\documentclass[12pt,a4paper]{article}
\usepackage[utf8]{inputenc}
\usepackage[T1]{fontenc}
\usepackage{amsmath,amssymb,amsfonts}
\usepackage{mathtools}
\usepackage{booktabs}
\usepackage{array}
\usepackage{geometry}
\usepackage{hyperref}
\usepackage{cite}
\geometry{margin=2.5cm}

\newcommand{\MeV}{\,\text{MeV}}
\newcommand{\GeV}{\,\text{GeV}}
\newcommand{\fpi}{f_\pi}
\newcommand{\qqbar}{\langle\bar{q}q\rangle}
\newcommand{\zch}{\mathfrak{z}}

\title{Three families from SUSY confinement: spectrum and mass predictions\\
of an $\mathcal{N}=1$ $\mathrm{SU}(3)\times\mathrm{SU}(2)$ theory}

\author{A.~Rivero\thanks{Universidad de Zaragoza.
Manuscript prepared with AI assistance (Claude/Anthropic).
Continues the program of \cite{Rivero2005}.}}
\date{March 2026}

\begin{document}
\maketitle

\begin{abstract}
We construct an $\mathcal{N}=1$ supersymmetric theory with gauge group
$\mathrm{SU}(3) \times \mathrm{SU}(2)$ and show that it uniquely selects
three colors and three generations through a Diophantine bootstrap.
Requiring that the scalar diquarks and mesons of an $\mathrm{SU}(N)$
confining theory fill supermultiplets imposes three interlocking
constraints --- $rs = 2N$, $r(r+1)/2 = 2N$, $r^2 + s^2 - 1 = 4N$ ---
whose only positive-integer solution is $N=3$, $r=3$, $s=2$. The five
light quarks carry an $\mathrm{SU}(5)$ flavor symmetry under which the
composites fill $\mathbf{24}\oplus\mathbf{15}\oplus\overline{\mathbf{15}}$,
a representation that is anomaly-free and asymptotically free.
The quark sector is described by $\mathrm{SU}(3)$ SQCD at
$N_f = N_c = 3$ in the Intriligator--Seiberg--Shih magnetic dual; the
lepton sector by $s$-confining $\mathrm{SU}(2)$ SQCD with $N_f = 3$.
Both sectors admit an O'Raifeartaigh deformation whose
Coleman--Weinberg-stabilized vacuum at $gv/m = \sqrt{3}$ fixes a
structural mass ratio $(2+\sqrt{3})^2 \approx 13.93$ and enforces an
algebraic identity $\langle z_k^2 \rangle = z_0^2$ on the
mass charges $m_k = (z_0 + z_k)^2$. From the strange quark mass alone,
the Lagrangian predicts $m_c = 1301$ MeV (+2.0\%), $m_b = 4233$ MeV
(+1.3\%), and $m_t = 171\,250$ MeV ($-$0.9\%). The lepton sector predicts
$m_\tau = 1776.97$ MeV (+0.006\%). In total five of the Standard Model's
thirteen charged-fermion-sector parameters are determined, reducing the
free parameter count from thirteen to eight.
\end{abstract}

%======================================================================
\section{Introduction}
%======================================================================

Of the roughly two dozen free parameters in the Standard Model, thirteen
belong to the charged fermion sector: nine masses and four CKM mixing
angles. No established principle relates them. The muon is 207 times
heavier than the electron; the top quark outweighs the up quark by a
factor of $8 \times 10^4$. These hierarchies are accommodated by the
Yukawa couplings, but the couplings themselves are inserted by hand. The
flavor problem --- why three generations, and why these masses --- remains
the central open question in particle physics.

This paper proposes a concrete $\mathcal{N}=1$ supersymmetric answer.
Its foundational relation is
\begin{equation}
m_\pi^2 - m_\mu^2 \;\approx\; \fpi^2\,,
\end{equation}
identifying the pion decay constant $\fpi \approx 92.2\MeV$ as the
SUSY-breaking order parameter: numerically,
$m_\pi^2 - m_\mu^2 = 8{,}316\MeV^2$ versus
$\fpi^2 = 8{,}501\MeV^2$ ($2.2\%$ discrepancy). If the pion and the
muon are superpartners split by soft SUSY breaking, the splitting
scale is~$\fpi$.

The construction begins from the observation that if the scalar partners
of quarks are composite (diquarks and mesons of a confining gauge
theory), then the requirement that they fill common supermultiplets is
not automatically satisfiable. The number of scalar diquarks formed from
$r$ flavors is $r(r+1)/2$; the number of scalar mesons from $r$ flavors
and $s$ antiflavors is $rs$; and these must jointly account for $4N$
bosonic degrees of freedom in a theory with $N$ colors. The resulting
system of Diophantine equations ---
$$rs = 2N, \qquad \frac{r(r+1)}{2} = 2N, \qquad r^2 + s^2 - 1 = 4N$$
--- admits a unique positive-integer solution: $N=3$, $r=3$, $s=2$. The
number of colors, the number of generations (encoded as $r + s - 1 = 4$,
giving four light quark flavors in addition to the strange), and the
representation content of the composites are all fixed by arithmetic.
This is the bootstrap: supersymmetry and confinement together select the
Standard Model's gauge and family structure.

The five light quarks $(u, d, s, c, b)$ then sit in the fundamental of
an $\mathrm{SU}(5)$ flavor symmetry. Under $\mathrm{SU}(5)_{\rm
flavor} \times \mathrm{SU}(3)_{\rm color}$, the composites --- mesons
in the adjoint-plus-singlet and diquarks in the symmetric tensor ---
fill the representations $\mathbf{24} \oplus \mathbf{15} \oplus
\overline{\mathbf{15}}$. This combination is anomaly-free and
asymptotically free with one-loop coefficient $b_0 = 31/3$. The top
quark, with Yukawa coupling $y_t \approx 1$, is external to the
confining dynamics; it does not form bound states.

The confining dynamics of the quark sector is that of $\mathrm{SU}(3)$
SQCD with $N_f = N_c = 3$ flavors, a theory whose infrared physics was
solved by Seiberg~\cite{Seiberg1995}. In the magnetic dual description, the degrees of
freedom are gauge-invariant mesons $M^i{}_j$, baryons $B$ and
$\tilde{B}$, and a Lagrange-multiplier singlet $X$ that enforces the
quantum-modified constraint $\det M - B\tilde{B} = \Lambda^6$. The
Seiberg seesaw gives meson vacuum expectation values $\langle M^j{}_j
\rangle = C/m_j$, inversely proportional to the quark masses. Adding an
O'Raifeartaigh deformation --- the minimal coupling $W \supset fX + m
\phi_a \phi_b + gX\phi_a^2$ --- breaks supersymmetry at a metastable
vacuum. The Coleman--Weinberg potential for the pseudo-modulus $X$,
combined with a non-canonical K\"ahler correction $\delta K = -|X|^4
/(12\mu^2)$, stabilizes the vacuum expectation value at $g\langle X
\rangle / m = \sqrt{3}$.

This stabilization has a remarkable consequence. The fermion mass matrix
at the O'Raifeartaigh vacuum has eigenvalues $(0,\; m_-,\; m_+)$ with
$m_\pm = (\sqrt{t^2 + 1} \pm t)\,m$ and $t = gv/m$. At $t = \sqrt{3}$
the eigenvalues become $(0,\; (2-\sqrt{3})\,m,\; (2+\sqrt{3})\,m)$,
whose ratio $m_+/m_- = (2+\sqrt{3})^2 \approx 13.93$ is a structural
constant independent of all free parameters. This is the
\emph{O'Raifeartaigh mass ratio}. Applied to the quark seed with
$m = m_s = 93.4$ MeV, it predicts $m_c = 13.93 \times m_s = 1301$ MeV,
to be compared with the PDG value $1275 \pm 25$ MeV (2.0\% deviation,
$1.0\sigma$).

The same vacuum enforces an algebraic identity on the mass charges.
Writing each mass as $m_k = \mathfrak{z}_k^2$ with
$\mathfrak{z}_k = z_0 + z_k$ and $\sum z_k = 0$, the eigenvalue
triple $(0, m_-, m_+)$ at $t = \sqrt{3}$ satisfies
$\langle z_k^2 \rangle = z_0^2$ --- the variance of the charge
perturbations equals the squared mean charge. This is not an empirical
observation extracted from data; it is a theorem about the
O'Raifeartaigh superpotential at its stabilized vacuum. This identity
serves as the organizing principle for all subsequent mass predictions.

To pass from the charm quark to the bottom quark, a second mechanism is
needed. On $\mathbb{R}^3 \times S^1$, the SU(3) gauge theory admits
magnetic bions --- correlated monopole--antimonopole pairs associated
with the three roots of the extended Dynkin diagram. Their contribution
to the effective K\"ahler potential modifies the mass eigenvalue relation
to $\sqrt{m_b} = 3\sqrt{m_s} + \sqrt{m_c}$, where the integer
coefficient is the number of colors. We call this the \emph{bion mass
relation}. Using the O'Raifeartaigh-predicted $m_c$, it gives
$m_b = 4233$ MeV (PDG: $4180 \pm 30$ MeV, 1.3\% deviation). The
mass-charge identity is approximately preserved:
$\langle z_k^2 \rangle / z_0^2 = 1.018$, deviating from unity by 0.9\%.

The chain extends to the top quark through the \emph{Yukawa eigenvalue
constraint}. At $\tan\beta = 1$, the SUSY Yukawa structure requires the
heavy triple $(m_c, m_b, m_t)$ to satisfy $\langle z_k^2 \rangle = z_0^2$
with all positive mass charges. This reduces to a quadratic equation for $\sqrt{m_t}$
whose physical solution, using the predicted $m_c$ and $m_b$ as inputs,
gives $m_t = 171\,250$ MeV (PDG: $172\,760 \pm 300$ MeV, 0.9\%
deviation). Thus from the single input $m_s$, the theory produces the
masses of all three heavy quarks.

The lepton sector is governed by the same algebraic structure. The
$\mathrm{SU}(2)$ confining theory with $N_f = 3$ is $s$-confining
($N_f = N_c + 1$); its infrared dynamics is captured by three
antisymmetric composites and a singlet, deformed by an O'Raifeartaigh
superpotential stabilized at the same $gv/m = \sqrt{3}$. Because the
SU(2) theory has no magnetic bion mechanism, the physical spectrum stays
close to the O'Raifeartaigh seed. The mass-charge identity
$\langle z_k^2 \rangle = z_0^2$ applied to $(m_e, m_\mu, m_\tau)$
directly predicts $m_\tau = 1776.97$ MeV
from the inputs $m_e$ and $m_\mu$, matching the PDG value $1776.86 \pm
0.12$ MeV to 0.006\%.

Two further predictions follow from the Lagrangian. The identification
of the Seiberg Lagrange multiplier $X$ with the NMSSM singlet $S$
generates a Higgs quartic through the $|F_X|^2$ term. At $\tan\beta =
1$ the tree-level Higgs mass is $m_h = \lambda v/\sqrt{2}$, and
requiring $m_h = 125.35$ GeV fixes $\lambda = 0.72$. The Cabibbo angle
emerges from the tachyonic off-diagonal meson condensate: the ratio of
condensate amplitudes in the $(u,s)$ and $(u,d)$ sectors is exactly
$\sqrt{m_d/m_s}$, reproducing the Gatto--Sartori--Tonin relation
$\sin\theta_C = \sqrt{m_d/m_s} = 0.224$ (PDG: 0.225).

The parameter budget of the theory is as follows. The Standard Model's
charged fermion sector has thirteen free parameters: nine masses and four
CKM angles. In this theory, the O'Raifeartaigh mass ratio determines
$m_c$ from $m_s$; the bion mass relation determines $m_b$ from $m_s$
and $m_c$; the Yukawa eigenvalue constraint determines $m_t$ from $m_c$
and $m_b$; the mass-charge identity determines $m_\tau$ from $m_e$
and $m_\mu$; and the Gatto--Sartori--Tonin relation determines
$\theta_{12}$ from $m_d$ and $m_s$. Five parameters are eliminated. The
remaining eight free parameters are $m_u$, $m_d$, $m_s$, $m_e$, $m_\mu$,
$\theta_{23}$, $\theta_{13}$, and $\delta_{\rm CP}$.

The paper is organized as follows. Section~\ref{sec:bootstrap} presents
the Diophantine bootstrap and proves the uniqueness of $N=3$, $r=3$,
$s=2$. Section~\ref{sec:diquarks} describes the $\mathrm{SU}(5)$ flavor
structure, the composite representations, and the anomaly cancellation.
Section~\ref{sec:lagrangian} constructs the $\mathrm{SU}(3)$ SQCD
Lagrangian in the Seiberg magnetic dual, derives the O'Raifeartaigh
vacuum, and proves the mass-charge identity $\langle z_k^2\rangle = z_0^2$.
Section~\ref{sec:predictions} derives the bion mass relation and the
Yukawa eigenvalue constraint, assembling the complete $m_s$ chain that
predicts $m_c$, $m_b$, and $m_t$. Section~\ref{sec:leptons} presents the
lepton sector as $s$-confining SU(2) SQCD with the same O'Raifeartaigh
structure. Section~\ref{sec:electroweak} discusses the NMSSM Higgs
sector, the Cabibbo angle, and the electroweak predictions.
Section~\ref{sec:spectrum} presents the complete particle spectrum and
parameter budget. Section~\ref{sec:open} discusses the open problems ---
the 15 versus $\overline{10}$ Pauli conflict for diquarks, the
nonperturbative origin of the bloom rotation, the mesino mass hierarchy,
and the embedding in $\mathrm{SO}(32)$ --- and Section~\ref{sec:conclusions}
concludes.


%======================================================================
\section{Bootstrap uniqueness}
\label{sec:bootstrap}
%======================================================================

The sBootstrap rests on a single postulate: every scalar of the
Supersymmetric Standard Model is a composite of two quarks (diquark,
color triplet) or of a quark--antiquark pair (meson, color singlet),
with the electric charge given by the sum of the constituents' charges.
We now show that matching the number of such composites to the
fermionic degrees of freedom of $N$ generations uniquely forces
$N = 3$.

\subsection{Setup and Diophantine system}
Let there be $N$ generations, each containing quarks of two types:
\emph{down} (charge $-1/3$) and \emph{up} (charge $+2/3$).  Not
every quark need participate as a string endpoint; we allow
$r$ light down-type quarks and $s$ light up-type quarks that form
composites.  The Standard Model with right-handed neutrinos has
$4N$ real (Weyl) fermionic degrees of freedom per charge sector in
each color channel.  Matching composites to fermions, charge sector
by charge sector, yields the following system.

\begin{enumerate}
\item \textbf{Charge $+1/3$ composites} (one up $\times$ one down,
analogous to charged sleptons): $rs$ complex scalars must match
$2N$ real fermionic d.o.f., giving
\begin{equation}
rs = 2N\,.
\end{equation}

\item \textbf{Charge $-2/3$ composites} (symmetric down--down pairs,
color sextet): $\binom{r+1}{2} = r(r+1)/2$ complex scalars must
match $2N$ real fermionic d.o.f., giving
\begin{equation}
\frac{r(r+1)}{2} = 2N\,.
\end{equation}

\item \textbf{Neutral composites} (classified by $SU(r+s)$, minus
the overall singlet trace; includes the partner of the right-handed
neutrino): $(r+s)^2 - 1 = r^2 + s^2 + 2rs - 1$ states, of which
$2rs$ are charged and already counted.  The neutral count is therefore
$r^2 + s^2 - 1$, to be matched to $4N$ real d.o.f., giving
\begin{equation}
r^2 + s^2 - 1 = 4N\,.
\end{equation}
\end{enumerate}

\noindent
The alternative neutral equation $r^2 + s^2 - 1 = 2N$ (left-handed
neutrinos only) produces no positive-integer solution when combined
with the first two; thus the sBootstrap \emph{requires}
right-handed neutrinos.

\subsection{The uniqueness theorem}

\medskip
\noindent\textbf{Theorem 1} (Generation uniqueness).
\textit{The system has a unique
solution over the positive integers:}
\begin{equation}
\boxed{(N,\,r,\,s) = (3,\,3,\,2)\,.}
\end{equation}

\medskip
\noindent\textit{Proof.}
Equating the right-hand sides of the first two equations
gives $rs = r(r+1)/2$, hence
\begin{equation}
s = \frac{r+1}{2}\,,
\end{equation}
which requires $r$ odd.  Substituting into the third equation:
\[
r^2 + \Bigl(\frac{r+1}{2}\Bigr)^{\!2} - 1 \;=\; 4\,\cdot\,\frac{r\,(r+1)/2}{2}
\;=\; r(r+1)\,.
\]
Expanding the left side,
\[
r^2 + \frac{r^2 + 2r + 1}{4} - 1 = r^2 + r\,,
\]
and clearing the denominator:
\[
4r^2 + r^2 + 2r + 1 - 4 = 4r^2 + 4r
\quad\Longrightarrow\quad
r^2 - 2r - 3 = 0\,,
\]
i.e.\ $(r-3)(r+1) = 0$.  Since $r > 0$, we obtain $r = 3$, whence
$s = 2$ and $N = rs/2 = 3$.
\hfill$\square$

\medskip
The solution has $r + s = 5$ light quarks acting as string endpoints,
furnishing the fundamental representation of an $SU(5)$ flavor group.
The top quark is the unique ``elephant''---too heavy to serve as a
constituent of the scalar composites---in precise agreement with the
Standard Model, where $m_t \gg \Lambda_{\text{QCD}}$ prevents
top-flavored hadrons from forming.

\subsection{Partial solutions and the hexagonal-number family}
The first two equations alone, without the
neutral constraint, require $2N$ to be a hexagonal number
$h_k = k(2k-1)$:
\[
2N = 1,\; 6,\; 15,\; 28,\; 45,\; 66,\; \ldots
\]
Of these, only $2N = 1$ and $2N = 6$ have $r \leq 16$ (the
asymptotic-freedom bound $2n_f < 33$ limits $r + s \leq 16$).
The first is inadmissible ($N = 1/2$); the second gives $N = 3$.

\subsection{Identification with $SU(5)$ representations}
With $r = 3$ and $s = 2$, the $r + s = 5$ light quarks form the
fundamental $\mathbf{5}$ of $SU(5)_f$.  The composites organize
into three $SU(5)$ representations:

\paragraph{Mesons (quark--antiquark, color singlet):}
The $\mathbf{5} \otimes \bar{\mathbf{5}}$ decomposes as
\begin{equation}
\mathbf{5} \otimes \bar{\mathbf{5}}
= \mathbf{24} \oplus \mathbf{1}\,.
\end{equation}
The $\mathbf{24}$ carries the sleptons (charged and neutral) and
the meson matrix.  Under
$SU(5) \supset SU(3)_d \times SU(2)_u$:
\begin{equation}
\mathbf{24} = (\mathbf{1},\mathbf{1})
+ (\mathbf{1},\mathbf{3})
+ (\mathbf{3},\mathbf{2})
+ (\bar{\mathbf{3}},\mathbf{2})
+ (\mathbf{8},\mathbf{1})\,.
\end{equation}

\paragraph{Diquarks (quark--quark, color triplet):}
The symmetric product
\begin{equation}
\mathbf{5} \otimes_S \mathbf{5}
= \mathbf{15}
\end{equation}
decomposes as
\begin{equation}
\mathbf{15} = (\mathbf{1},\mathbf{3})
+ (\mathbf{3},\mathbf{2})
+ (\mathbf{6},\mathbf{1})\,,
\end{equation}
and its conjugate $\overline{\mathbf{15}}$ provides the
antidiquarks.

\medskip
\noindent\textbf{Proposition 1} (Symmetric product is forced).
\textit{The counting equation selects the
symmetric product $\mathbf{15}$, not the antisymmetric
$\overline{\mathbf{10}}$.}

\medskip
\noindent\textit{Proof.}
The charge $-2/3$ composites are down--down pairs.
Under $SU(3)_d$, the symmetric product
$\mathbf{3} \otimes_S \mathbf{3} = \mathbf{6}$ has dimension
$r(r+1)/2 = 6$, matching the counting equation.  The antisymmetric
product $\mathbf{3} \otimes_A \mathbf{3} = \bar{\mathbf{3}}$ has
dimension $r(r-1)/2 = 3 \neq 2N = 6$.  The symmetric color sextet,
and hence the $\mathbf{15}$, is the unique choice consistent
with the sBootstrap counting.
\hfill$\square$

\subsection{Charge census of $\mathbf{24} \oplus \mathbf{15}
\oplus \overline{\mathbf{15}}$}
Treating the $SU(5)$ simultaneously as a Georgi--Glashow
unification group~\cite{Rivero2005}, one assigns standard electric
charges via $Q_{\text{em}} = T_3 - Y/6$ to every component.
Table~\ref{tab:census} displays the resulting charge census for
$\mathbf{24} \oplus \mathbf{15} \oplus \overline{\mathbf{15}}$.

\begin{table}[ht]
\centering
\begin{tabular}{lcccccc@{\qquad}c}
\toprule
& \multicolumn{6}{c}{Number of states at $|Q_{\text{em}}|$} & \\
\cmidrule(lr){2-7}
& $0$ & $\tfrac{1}{3}$ & $\tfrac{2}{3}$ & $1$
& $\tfrac{4}{3}$ & $2$ & Dim \\
\midrule
$\mathbf{24}$ & 10 & $2\!\times\!3$ & ---
& $2\!\times\!1$ & $2\!\times\!3$ & --- & 24 \\[2pt]
$\mathbf{15}$ & 1 & 3 & 9
& 1 & --- & 1 & 15 \\[2pt]
$\overline{\mathbf{15}}$ & 1 & 3 & 9
& 1 & --- & 1 & 15 \\[2pt]
\midrule
Total & 12 & $2\!\times\!6$ & $2\!\times\!9$
& $2\!\times\!2$ & $2\!\times\!3$ & $2\!\times\!1$ & 54 \\
\bottomrule
\end{tabular}
\caption{Charge census of
$\mathbf{24} \oplus \mathbf{15} \oplus \overline{\mathbf{15}}$
of $SU(5)$.
Notation ``$2\!\times\!k$'' means $k$ states at $+|Q_{\text{em}}|$ and $k$
at $-|Q_{\text{em}}|$; the $\mathbf{24}$ is self-conjugate.
Charges $\pm 4/3$ in the $\mathbf{24}$ arise from the leptoquark
components $(\mathbf{3},\mathbf{2})_{-5/6}$ and
$(\bar{\mathbf{3}},\mathbf{2})_{5/6}$.}
\label{tab:census}
\end{table}

\noindent
Three cross-checks tie the census to the Diophantine system:
\begin{itemize}
\item \textbf{Neutral composites:} The 12 states at $Q_{\text{em}} = 0$
match $r^2 + s^2 - 1 = 9 + 4 - 1 = 12 = 4N$.
\item \textbf{Charge $+1/3$ composites:}
$rs = 3 \times 2 = 6$ composites, matching
$rs = 6 = 2N$.
\item \textbf{Charge $-2/3$ composites:}
$r(r+1)/2 = 6$ symmetric pairs fill the color sextet
$(\mathbf{6},\mathbf{1})_{-2/3} \subset \mathbf{15}$.
\end{itemize}

\subsection{Anomaly cancellation and asymptotic freedom}
The representation $\mathbf{24} \oplus \mathbf{15} \oplus
\overline{\mathbf{15}}$ must be anomaly-free if it is to appear
in a consistent gauge theory.

\medskip
\noindent\textbf{Proposition 2} (Anomaly cancellation).
\textit{The cubic $SU(5)$ anomaly of $\mathbf{24} \oplus
\mathbf{15} \oplus \overline{\mathbf{15}}$ vanishes.}

\medskip
\noindent\textit{Proof.}
The anomaly coefficient $A(\mathbf{R})$ for $SU(n)$ satisfies
$A(\text{adjoint}) = 0$ (the adjoint is real),
$A(\mathbf{15}) = n + 4 = 9$, and
$A(\overline{\mathbf{15}}) = -(n+4) = -9$.  Therefore
\[
A(\mathbf{24}) + A(\mathbf{15}) + A(\overline{\mathbf{15}})
= 0 + 9 + (-9) = 0\,.
\]
\hfill$\square$

\medskip
\noindent\textbf{Proposition 3} (Asymptotic freedom).
\textit{An $SU(5)$ gauge theory with matter in the
$\mathbf{24} \oplus \mathbf{15} \oplus \overline{\mathbf{15}}$
is asymptotically free.}

\medskip
\noindent\textit{Proof.}
The one-loop beta-function coefficient is
\[
b_0 = \frac{11}{3}\,C_2(G) - \frac{2}{3}\,\sum_{\text{Weyl}} T(\mathbf{R})\,,
\]
where $C_2(G) = n = 5$ for $SU(5)$, $T(\mathbf{adj}) = n = 5$,
and $T(\mathbf{15}) = (n+2)/2 = 7/2$.  The Weyl fermion content
consists of one adjoint and one $\mathbf{15} + \overline{\mathbf{15}}$,
giving
$\sum T = 5 + 7/2 + 7/2 = 12$.  Therefore
\[
b_0 = \frac{11}{3} \times 5 - \frac{2}{3} \times 12
= \frac{55 - 24}{3} = \frac{31}{3} > 0\,.
\]
\hfill$\square$

\subsection{Summary of the counting argument}
The logical chain is:
\begin{enumerate}
\item The sBootstrap postulate (composites = scalars of the SSM)
yields three Diophantine equations in $(N, r, s)$.
\item The unique positive-integer solution is $(3, 3, 2)$:
three generations, with the symmetric $\mathbf{15}$ forced over
the antisymmetric $\overline{\mathbf{10}}$.
\item The solution organizes into $SU(5)$ representations
$\mathbf{24} \oplus \mathbf{15} \oplus \overline{\mathbf{15}}$,
which are anomaly-free and asymptotically free.
\item Right-handed neutrinos are required: without them, no
integer solution exists.
\item The top quark is excluded from the composite sector by
the counting itself ($s = 2 < N = 3$), matching its
kinematic decoupling ($m_t \gg \Lambda_{\text{QCD}}$).
\end{enumerate}


%======================================================================
\section{SU(5) flavor and the diquark sector}
\label{sec:diquarks}
%======================================================================

\subsection{Counting partners}

In the full sBootstrap~\cite{Rivero2005}, the scalar partners of
the Standard Model fermions include both mesons ($q\bar{q}$) and
diquarks ($qq$). For five active flavors, the meson matrix gives
$5 \times 5 = 25$ scalars. The diquarks $[q_i q_j]$ form a
representation of $SU(5)_f$:
\begin{equation}
\mathbf{5} \otimes \mathbf{5} = \overline{\mathbf{10}}_A
\oplus \mathbf{15}_S\,.
\end{equation}

The symmetric $\mathbf{15}_S$ gives 15 diquarks including
same-flavor states $[uu]$, $[dd]$, \ldots, while the antisymmetric
$\overline{\mathbf{10}}_A$ gives 10 diquarks with $i \neq j$ only.

\subsection{The symmetric representation}

The sBootstrap SUSY generator acts asymmetrically on the two
fermion sectors. When it acts on a \emph{lepton}, it produces a
meson: $\mathcal{Q} \cdot \ell \sim q\bar{q}$. Since the quark and antiquark
are distinguishable, the meson lies in the full
$\mathbf{5} \otimes \bar{\mathbf{5}} = \mathbf{24} + \mathbf{1}$.
There is no Pauli constraint, and the meson spectrum is perfectly
reproduced.

When it acts on a \emph{quark}, it produces a diquark:
$\mathcal{Q} \cdot q_i \sim q_i q_j$. The bosonic SUSY generator applied to
a single flavor index $i$ naturally produces the \textbf{symmetric}
product---the $\mathbf{15}_S$ representation, not the antisymmetric
$\overline{\mathbf{10}}$. This is the representation identified
in the $SO(32)$ decomposition~\cite{Rivero2005}: the adjoint of
$SO(32)$ under $SU(5)_f$ decomposes to include the
$\mathbf{15} + \overline{\mathbf{15}}$ (diquarks) alongside the
$\mathbf{1} + \mathbf{24}$ (mesons).

\subsection{The Pauli theorem for scalar diquarks}

The requirement of the symmetric $\mathbf{15}$ conflicts with
Fermi statistics---not merely for $L=0$ diquarks, but for
\emph{all} scalar diquarks in the color $\bar{3}$ channel,
regardless of internal orbital structure.

\medskip
\noindent\textbf{Theorem.}
\textit{For a spin-0 ($J=0$) diquark in the color $\bar{3}$
(antisymmetric) representation, the flavor wave function is
antisymmetric for \textbf{all} values of the orbital angular
momentum $L$ and spin $S$.}

\medskip
\noindent\textit{Proof.}
$J = 0$ requires $L = S$ (since $|L - S| \leq J \leq L + S$).
The exchange symmetry of the spin-orbital wave function is
$(-1)^{L+S+1} = (-1)^{2L+1} = -1$ for all $L$. Combined with
the color $\bar{3}$ factor $(-1)$, the non-flavor product is
$(-1)(-1) = +1$ (symmetric). Total antisymmetry then requires
flavor to be antisymmetric: $\overline{\mathbf{10}}$. \hfill$\square$

\medskip
This result is purely kinematic---independent of dynamics,
binding mechanisms, or spatial extent. The key cases:

\begin{center}
\begin{tabular}{ccccc}
\toprule
$L$ & $S$ & Spin-orbital $(-1)^{2L+1}$ & $\times$ Color & Flavor \\
\midrule
0 & 0 & $-1$ & $\times\,(-1) = +1$ & must be A \\
1 & 1 & $-1$ & $\times\,(-1) = +1$ & must be A \\
2 & 2 & $-1$ & $\times\,(-1) = +1$ & must be A \\
\bottomrule
\end{tabular}
\end{center}

The theorem has a sharp consequence: \textbf{the sBootstrap
diquark in the $\mathbf{15}$ cannot be a $qq$ composite.} It
must be either:
\begin{enumerate}
\item An algebraically generated object---the image of the
Miyazawa supercharge~\cite{Miyazawa, CattoGursey} acting on
an antiquark, whose symmetry is determined by the superalgebra
rather than by the statistics of constituent quarks; or
\item A fundamental string-theoretic object, as in the $SO(32)$
decomposition~\cite{Rivero2005}, where the $\mathbf{15}$
arises from the adjoint representation of the gauge group and
the open-string diquark is not built from point-like quarks.
\end{enumerate}

The contrast with the meson sector is instructive. When the
SUSY generator acts on a lepton, it produces a meson
($q\bar{q}$): the quark and antiquark are distinguishable,
no Pauli constraint applies, and the meson spectrum is
reproduced without tension. When it acts on a quark, it
produces a diquark in the ``wrong'' representation---and the
theorem shows this cannot be fixed by any internal structure
of a $qq$ pair. This is an unresolved obstruction: the sBootstrap requires
the symmetric $\mathbf{15}$, but QCD composite diquarks cannot
fill it.


%======================================================================
\section{The SUSY Lagrangian}
\label{sec:lagrangian}
%======================================================================

\subsection{The Seiberg framework}

The natural arena is $\mathcal{N}=1$ SQCD with gauge group $SU(3)_c$
and $N_f$ flavors of chiral superfields $Q^i$, $\bar{Q}_j$ in the
fundamental and antifundamental representations. For $N_f = N_c = 3$
(the case relevant to light QCD), Seiberg's exact results
\cite{Seiberg1995} establish that the theory confines, with the
low-energy degrees of freedom being the composite meson superfield
\begin{equation}
M^i{}_j = Q^i \cdot \bar{Q}_j\,,
\end{equation}
a $3\times 3$ matrix of chiral superfields, together with baryons
$B = \epsilon_{ijk}Q^iQ^jQ^k$ and $\bar{B}$. These satisfy the
quantum-modified constraint
\begin{equation}
\det M - B\bar{B} = \Lambda^{6}\,,
\end{equation}
where $\Lambda$ is the dynamical scale.

Each entry $M^i{}_j$ is a chiral superfield containing a complex
scalar (the meson) and a Weyl fermion (the lepton, in the sBootstrap
identification). The scalar and fermion are degenerate in the SUSY
limit; SUSY breaking splits them.

\subsection{F-term breaking: the O'Raifeartaigh mechanism}

The O'Raifeartaigh model~\cite{ORaifeartaigh} is the simplest
realization of spontaneous $F$-term SUSY breaking. With three chiral
superfields, the superpotential
\begin{equation}
W = f\,\Phi_0 + m\,\Phi_1\Phi_2 + g\,\Phi_0\Phi_1^2
\end{equation}
gives $F$-terms:
\begin{align}
F_{\Phi_0}^* &= f + g\,\phi_1^2\,, \\
F_{\Phi_1}^* &= m\,\phi_2 + 2g\,\phi_0\phi_1\,, \\
F_{\Phi_2}^* &= m\,\phi_1\,.
\end{align}
If $f \neq 0$ and $m \neq 0$, then $F_{\Phi_2}^* = 0$ requires
$\phi_1 = 0$, but then $F_{\Phi_0}^* = f \neq 0$: SUSY is
\textbf{necessarily broken}. The key features are:
\begin{enumerate}
\item A massless field ($\Phi_0$) exists at tree level---it is the
pseudo-modulus.
\item The breaking is triggered because $F_{\Phi_0}$ cannot vanish.
\item After including loop corrections, $\Phi_0$ acquires a mass
(the Coleman--Weinberg potential lifts the flat direction).
\end{enumerate}

The fermionic mass matrix at the vacuum ($\phi_1 = 0$,
$\phi_0 = v$) is
\begin{equation}
  W_{ij} = \begin{pmatrix} 0 & 0 & 0 \\ 0 & 2gv & m \\ 0 & m & 0 \end{pmatrix},
\end{equation}
with eigenvalues $0$ (the goldstino) and $m_\pm = \sqrt{g^2v^2 + m^2} \pm gv$.
Setting the mass ratio to the seed value
$\sqrt{m_-}/\sqrt{m_+} = 2 - \sqrt{3}$
requires $\sqrt{t^2+1} - t = 2 - \sqrt{3}$, where $t = gv/m$,
which gives the unique solution
\begin{equation}
  \frac{gv}{m} = \sqrt{3}\,.
\end{equation}
At this vacuum, $g^2v^2 + m^2 = 4m^2$, so
$m_\pm = (2 \pm \sqrt{3})\,m$.

\subsection{Abelian charges and the mass-charge identity}

Write each fermion mass as the square of an abelian charge,
\begin{equation}\label{eq:mass-charge-def}
  m_k = \zch_k^2, \qquad \zch_k = z_0 + z_k, \qquad k=1,2,3,
\end{equation}
with $z_0 \equiv \bar{\zch} = \tfrac{1}{3}\sum_k \zch_k$ and the tracelessness
constraint $\sum_k z_k = 0$.
Then the mass-charge identity $\langle z_k^2\rangle = z_0^2$ takes the form of a
\emph{charge-variance identity}: expanding
\begin{equation}
  \frac{1}{3}\sum_k z_k^2
  = \frac{1}{3}\sum_k \zch_k^2 - z_0^2
  = \frac{\sum m_k}{3} - \frac{(\sum\sqrt{m_k})^2}{9} ,
\end{equation}
so that
$\langle z_k^2 \rangle = z_0^2$ is equivalent to
$\tfrac{1}{3}\sum m_k = 2z_0^2$, which rearranges to
\begin{equation}
\boxed{
  \frac{\sum m_k}{\bigl(\sum\sqrt{m_k}\bigr)^2} = \frac{2}{3}
}\,.
\end{equation}
In statistical language the identity states that the population variance of
the charges equals the squared mean:
$\operatorname{Var}(\zch) = \bar{\zch}^{\,2}$.

\paragraph{Angular parametrisation.}
The constraint $\sum z_k = 0$ restricts $(z_1,z_2,z_3)$ to a
two-dimensional hyperplane in $\mathbb{R}^3$, parametrised as
\begin{equation}\label{eq:zk-angular}
  z_k = r\cos\!\bigl(\delta + \tfrac{2\pi k}{3}\bigr), \qquad k=1,2,3.
\end{equation}
For any $\delta$ the tracelessness $\sum\cos(\delta+2\pi k/3)=0$ is
automatic, and the identity
$\sum\cos^2(\delta+2\pi k/3) = 3/2$ gives
$\langle z_k^2\rangle = r^2/2$.
The charge-variance identity $\langle z_k^2\rangle = z_0^2$ then fixes
\begin{equation}
  r = \sqrt{2}\,z_0,
\end{equation}
and the angle $\delta$---the \emph{bloom angle}---parametrises the family
of mass triples at fixed vacuum charge $z_0$.
The full charges read
$\zch_k = z_0\bigl[1 + \sqrt{2}\cos(\delta + 2\pi k/3)\bigr]$.

\paragraph{The seed angle.}
At the O'Raifeartaigh seed, $m_1=0$ forces $\zch_1=0$, hence
$z_1 = -z_0$.  From~\eqref{eq:zk-angular},
\begin{equation}
  r\cos\delta = -z_0 = -\frac{r}{\sqrt{2}}
  \quad\Longrightarrow\quad
  \cos\delta = -\frac{1}{\sqrt{2}},
  \qquad \delta = \frac{3\pi}{4}.
\end{equation}
Every O'Raifeartaigh seed sits at $\delta=3\pi/4$.
The remaining two eigenvalues satisfy
$t^2 - 4mt + m^2 = 0$
with $z_0^2 = 2m/3$, giving $m_\pm = (2\pm\sqrt{3})\,m$, which is the
spectrum at $gv/m=\sqrt{3}$.

\paragraph{Determinant condition.}
Their product
$m_+m_- = m^2$ is the \emph{determinant condition}: it equals
$\det W_{ij}$ restricted to the massive $2\times 2$ block.
That this follows from $\langle z_k^2\rangle = z_0^2$
is direct: at the seed, $z_2+z_3=z_0$ and
$z_2 z_3 = -z_0^2/2$, so
\begin{equation}
  m_+ m_- = \bigl[(z_0{+}z_2)(z_0{+}z_3)\bigr]^2
  = \Bigl[z_0^2 + z_0^2 - \tfrac{z_0^2}{2}\Bigr]^2
  = \tfrac{9z_0^4}{4}
  = m^2.
\end{equation}
Conversely, requiring $m_+m_-=m^2$ with $m_++m_-=4m$ forces
$\langle z_k^2\rangle = z_0^2$.
Thus the mass-charge identity at the seed is the determinant of the
superpotential mass matrix.

\paragraph{Generality of the determinant condition.}
For a general O'Raifeartaigh superpotential
$W = f\Phi_0 + \frac{1}{2}m_{kl}\Phi_k\Phi_l
+ \frac{1}{2}\Phi_0\,g_{kl}\Phi_k\Phi_l$,
the massive fermion block is $M = m + v\,g$.
Then $\det M = \det m$ for all~$v$ if and only if $m^{-1}g$ is
\emph{nilpotent}. For $n=2$ this reduces to two conditions:
$\mathrm{Tr}(\mathrm{adj}(m)\cdot g) = 0$ and $\det g = 0$.
Both are satisfied automatically when $m$ is antidiagonal
and $g$ has a single diagonal entry (as in the standard model),
since $\mathrm{adj}(m)\cdot g$ is then strictly lower-triangular.
The product $m_+m_- = m^2$ is therefore not a numerical accident
at $gv/m = \sqrt{3}$; it holds for \emph{all} values of the
pseudo-modulus.

\paragraph{Signed charges and the quark sector.}
The charge $\zch_k$ need not be positive: a negative sign
$\zch_1 = -\sqrt{m_1}$ still gives $m_1 = \zch_1^2>0$
while changing the sign of the field's mass contribution in the
superpotential.

\paragraph{Vacuum-charge doubling.}
The seed triple $(0,\sqrt{m_s},\sqrt{m_c})$ has
$v_0^{\text{seed}} = (\sqrt{m_s}+\sqrt{m_c})/3$.
If $\sqrt{m_b} = 3\sqrt{m_s}+\sqrt{m_c}$, substitution gives
\begin{equation}\label{eq:v0-doubling}
  v_0^{\text{bloom}}
  = \frac{-\sqrt{m_s}+\sqrt{m_c}+3\sqrt{m_s}+\sqrt{m_c}}{3}
  = \frac{2(\sqrt{m_s}+\sqrt{m_c})}{3}
  = 2\,v_0^{\text{seed}}.
\end{equation}
Conversely, $v_0^{\text{bloom}}=2v_0^{\text{seed}}$ requires
$\sqrt{m_b}=3\sqrt{m_s}+\sqrt{m_c}$,
predicting $m_b = 4177$\,MeV (PDG: $4180\pm30$, $0.1\sigma$).
The bion relation and the vacuum-charge doubling are
algebraically equivalent statements: one is the other rewritten in the
charge language.

\medskip
The mass-charge identity is a theorem about the O'Raifeartaigh superpotential at
$gv/m = \sqrt{3}$, not an empirical observation.  The angular parametrisation,
determinant condition, and vacuum-charge doubling are its corollaries.

\subsection{K\"ahler stabilization of the pseudo-modulus}

In the minimal O'Raifeartaigh model with canonical K\"ahler potential,
the CW potential is monotonically increasing and selects
$\langle X\rangle = 0$.  A negative quartic K\"ahler correction
\begin{equation}
  K = |X|^2 + \frac{c}{\Lambda_K^2}\,|X|^4\,,
  \qquad c < 0
\end{equation}
modifies the tree-level potential to
$V_{\text{tree}} = f^2/(1 + 4c\,v^2/\Lambda_K^2)$, which
diverges at $v_{\text{pole}} = \Lambda_K/(2\sqrt{|c|})$
where the K\"ahler metric vanishes.  The field space terminates
at $v_{\text{pole}}$: the proper distance to the pole is finite,
$\sigma_{\max} = \sqrt{3}\pi m/(4g)$.

Setting $v_{\text{pole}} = \sqrt{3}\,m/g$ gives the exact condition
\begin{equation}
  c = -\,\frac{\Lambda_K^2}{12\,m^2}\,.
\end{equation}
With $\Lambda_K = m$ this is $c = -1/12$, an algebraic result:
$\Lambda_K/(2\sqrt{1/12}) = \sqrt{12}/2 = \sqrt{3}$.

\textbf{Open problem.} Both $V_{\text{tree}}$ and $V_{\text{CW}}$
are monotonically increasing functions of $v$: the one-loop
effective potential is minimized at $v = 0$, not at the pole.
The pole creates an impenetrable wall but does not attract.
Stabilizing the pseudo-modulus at $v = \sqrt{3}\,m/g$ requires
a nonperturbative negative contribution to the potential that
competes with the perturbative terms near the pole. The
three-instanton operator (Section~\ref{sec:open}) is a
candidate but has not been shown to produce the required
minimum.

\subsection{The O'Raifeartaigh structure in SQCD}

\paragraph{The quantum constraint.}
In Seiberg's SQCD with $N_f = N_c = 3$, the confined degrees of freedom
are the meson matrix $M^i{}_j = Q^i\bar Q_j/\Lambda$, the baryons
$B = \epsilon_{ijk}Q^iQ^jQ^k/\Lambda^3$ and
$\bar B = \epsilon^{ijk}\bar Q_i\bar Q_j\bar Q_k/\Lambda^3$,
subject to the quantum-modified constraint
$\det M - B\bar B = \Lambda^6$.
A Lagrange multiplier $X$ enforces this:
\begin{equation}
  W = X\bigl(\det M - B\bar B - \Lambda^6\bigr) + \sum_j m_j\, M^j{}_j\,.
\end{equation}

\paragraph{The Seiberg seesaw.}
Setting $B = \bar B = 0$ and diagonal $M = \mathrm{diag}(M_1,M_2,M_3)$,
the F-term equation $\partial W/\partial M^j{}_j = 0$ gives
\begin{equation}
  m_j + X\,\frac{\det M}{M_j} = 0
  \qquad(\text{no sum on } j)\,,
\end{equation}
so that $M_j = C/m_j$ with $C = -X\det M$ flavor-independent.
The constraint $\det M = \Lambda^6$ fixes
$C = \Lambda^2(m_1 m_2 m_3)^{1/3}$. Heavy quarks produce light
mesons: this is the Seiberg seesaw.

\paragraph{F-term SUSY breaking.}
Adding an O'Raifeartaigh deformation $\Delta W = f\,X + g\,X\,\phi^2/2$
(where $\phi$ is an off-diagonal meson fluctuation), the F-term
for $X$ becomes
\begin{equation}
  F_X^* = \bigl(\det M - B\bar B - \Lambda^6\bigr) + f
  + \tfrac{1}{2}g\,\phi^2\,.
\end{equation}
At the meta-stable vacuum where $\det M \approx \Lambda^6$ and
$\phi = 0$, the constraint term vanishes but $F_X^* = f \neq 0$.
Setting all F-terms to zero simultaneously would require
$\det M - \Lambda^6 = -f$, contradicting the remaining F-term
equations which drive $\phi \to 0$ and $\det M \to \Lambda^6$.
Supersymmetry is therefore spontaneously broken with vacuum
energy $V_0 = |f|^2$.

\paragraph{ISS parameter map.}
In the ISS magnetic dual with gauge group $SU(N_f - N_c)$,
the rank condition provides an independent argument:
$\tilde q \, q$ is an $N_f \times N_f$ matrix of rank
$\tilde N_c = N_f - N_c < N_f$, so it cannot equal
$\mu^2\,\mathbb{1}_{N_f}$. The $N_c$ frustrated F-terms
give $V_0 = N_c|h\mu^2|^2$.
The O'Raifeartaigh parameters map as
$g \to h$, $m \to h\mu$, $v \to \langle X\rangle$, so the
condition $gv/m = \sqrt{3}$ becomes
$\langle X\rangle/\mu = \sqrt{3}$---the magnetic Yukawa~$h$ cancels.

\subsection{The mass-charge identity is not an RG attractor}

One might hope that the identity
$\langle z_k^2\rangle = z_0^2$ is a dynamical attractor:
that three masses evolving under SUSY-breaking dynamics flow toward
the $\langle z_k^2\rangle = z_0^2$ surface. However, the one-loop soft-mass RG
equations in $\mathcal{N}=1$ SQCD with universal Casimirs and
negligible Yukawas give $dm_i/dt = \gamma(M_g^2 - m_i)$, which
drives all masses to a common value $m_i \to M_g^2$, yielding
$\langle z_k^2\rangle/z_0^2 \to 0$---the \emph{opposite} extreme.
The algebraic preservation condition on the
$\langle z_k^2\rangle = z_0^2$ surface is incompatible with any
uniform attractor flow (by the Cauchy--Schwarz inequality).

The identity therefore does \emph{not} arise from simple RG running.
It must instead be a \emph{boundary condition} on the
superpotential---as in the O'Raifeartaigh realization, where
$gv = \sqrt{3}\,m$ produces the exact seed. The condition is set at
the scale of SUSY breaking and then preserved by QCD-invariant mass
ratios.

\subsection{The full Lagrangian}

Consider $\mathcal{N}=1$ SQCD with gauge group $SU(3)_c$ and $N_f=5$ flavors
of chiral superfields $Q^i$, $\bar Q_j$ ($i,j=1\ldots 5$, corresponding to
$u,d,s,c,b$). The total Lagrangian is
\begin{equation}
  \mathcal{L}
  = \int\!d^4\theta\,K
  + \left[\int\!d^2\theta\,W + \text{h.c.}\right]
  + \mathcal{L}_D
  + \mathcal{L}_{\text{soft}}\,,
\end{equation}
with K\"ahler potential (canonical plus stabilization and bion corrections),
\begin{equation}
  K = \underbrace{\operatorname{Tr}(M^\dagger M) + |X|^2 + |B|^2
    + |\bar B|^2 + |H_u|^2 + |H_d|^2 + |S|^2}_{K_{\text{can}}}
  \underbrace{{}-\,\frac{|X|^4}{12\,\mu^2}}_{K_{\text{stab}}}
  + \underbrace{\frac{\zeta^2}{\Lambda^2}\,
    \Bigl|\sum_{k=1}^{N_c} s_k \sqrt{\hat m_k}\,\Bigr|^2}_{K_{\text{bion}}}\,,
\end{equation}
where the stabilization term pins the pseudo-modulus at
$\langle X\rangle = \sqrt{3}\,\mu$ via the K\"ahler pole,
$\zeta \propto e^{-S_0/(2N_c)}$ is the monopole fugacity,
$\hat m_k$ are the current quark masses, and
$s_k = \pm 1$ are monopole-instanton signs encoding the seed-to-bloom phase.

The superpotential combines mass, constraint, and Yukawa terms:
\begin{equation}
\begin{split}
  W &= \underbrace{\operatorname{Tr}(\hat m\, M)}_{\text{chiral mass}}
    + \underbrace{X\!\left(\det M^{(3)} - B\bar B
      - \Lambda_{\text{eff}}^{6}\right)}_{\text{Seiberg constraint}}
    + \underbrace{c_3\,\frac{(\det M^{(3)})^3}
      {\Lambda_{\text{eff}}^{18}}}_{\text{3-instanton}}\\
  &\quad  + \underbrace{y_c\, H_u\, M^c{}_c
      + y_b\, H_d\, M^b{}_b
      + y_t\, H_u\, Q^t \bar Q_t}_{\text{Yukawa}}
    + \underbrace{\lambda_S\, S\, H_u\!\cdot\!H_d
      + \tfrac{\kappa}{3}\,S^3}_{\text{NMSSM}}\,,
\end{split}
\end{equation}
where $M^{(3)}$ is the $3\times 3$ light-flavor block ($u,d,s$)
and $\Lambda_{\text{eff}}$ is the dynamical scale of the confined
$N_f = N_c = 3$ theory.

The $D$-term Lagrangian includes an optional Fayet--Iliopoulos
parameter $\xi$ for $U(1)_{\text{em}}$:
\begin{equation}
  \mathcal{L}_D
  = -\frac{e^2}{2}\!\left(\sum_i Q_i\,|\phi_i|^2
    + \xi\right)^{\!2}.
\end{equation}
Finally, soft SUSY-breaking terms arise from a spurion
$\langle F_S\rangle = \fpi^2$:
\begin{equation}
  \mathcal{L}_{\text{soft}}
  = -\tilde m^2 \operatorname{Tr}(M^\dagger M)
  - \bigl(B_\mu\operatorname{Tr}(\hat m\, M)
    + A_y\bigl(y_c\, H_u\, M^c{}_c + y_b\, H_d\, M^b{}_b\bigr) + \text{h.c.}\bigr)\,,
  \quad
  \tilde m^2 \sim \fpi^2\,.
\end{equation}
The soft scalar mass $\tilde m^2 \sim \fpi^2 \approx (92\MeV)^2$
generates the universal meson--lepton splitting
$m^2_{\text{meson}} - m^2_{\text{lepton}} = \fpi^2$.

\subsection{Scope of the Lagrangian}

This Lagrangian produces: (i)~the O'Raifeartaigh seed condition ($gv/m = \sqrt{3}$);
(ii)~the meson--lepton mass splitting ($\fpi^2$);
(iii)~the $v_0$-doubling condition via the bion K\"ahler
correction $K_{\text{bion}}$, whose minimum enforces
$\sqrt{m_b} = 3\sqrt{m_s} + \sqrt{m_c}$;
(iv)~the Yukawa couplings for the heavy sector;
(v)~the $\cos(9\delta)$ potential that selects
the lepton mass-spectrum phase, via the three-instanton operator
$(\det M^{(3)})^3/\Lambda_{\text{eff}}^{18}$.

The two-doublet structure has a direct consequence for the mass
triples. Since the quark mass triples obligatorily mix up-type
and down-type quarks, the translation from
masses to Yukawa couplings depends on~$\tan\beta$.
For the $(c,b,t)$ triple, with $c$ and $t$ receiving mass
from $H_u$ and $b$ from~$H_d$:
\begin{equation}
  \frac{\sum m_k}{(\sum\sqrt{m_k})^2}\bigg|_{c,b,t}
  = \frac{y_c^2\sin^2\!\beta + y_b^2\cos^2\!\beta + y_t^2\sin^2\!\beta}
  {(y_c\sin\beta + y_b\cos\beta + y_t\sin\beta)^2}\,,
\end{equation}
which satisfies $\langle z_k^2\rangle = z_0^2$ for the Yukawa eigenvalues
only when $\sin\beta = \cos\beta$,
i.e., $\tan\beta = 1$. Thus $\tan\beta = 1$ is the unique value
at which the mass-charge identity on \emph{masses} is equivalent
to the same identity on \emph{Yukawa couplings} for mixed
up/down-type triples.


%======================================================================
\section{Mass predictions from the superpotential}
\label{sec:predictions}
%======================================================================

The preceding sections established the Lagrangian: an $\mathcal{N}=1$
theory with K\"ahler potential $K = K_{\mathrm{can}} + K_{\mathrm{stab}} + K_{\mathrm{bion}}$,
superpotential $W$ containing the Seiberg constraint, chiral mass terms,
and Yukawa couplings, and soft breaking from a spurion with
$\langle F\rangle = \fpi^2$. We now extract the mass predictions.
The logic is sequential: each step uses the output of the previous one,
so the entire heavy-quark spectrum flows from a single input mass.

\subsection{The O'Raifeartaigh mass ratio}

The O'Raifeartaigh sector of the superpotential,
\begin{equation}
W_{\mathrm{OR}} = f\,\Phi_0 + m\,\Phi_1\Phi_2 + g\,\Phi_0\Phi_1^2\,,
\end{equation}
has a metastable SUSY-breaking vacuum at $\phi_1 = 0$, with
$\Phi_0$ a pseudo-modulus. At the vacuum
$gv/m = \sqrt{3}$, the fermion eigenvalues are
$m_\pm = (2 \pm \sqrt{3})\,m$,
and the mass ratio of the two nonzero fermion eigenvalues is
a structural constant:
\begin{equation}
R_{\mathrm{OR}}
= \frac{m_+}{m_-}
= \frac{2 + \sqrt{3}}{2 - \sqrt{3}}
= (2 + \sqrt{3})^2
= 7 + 4\sqrt{3}
\approx 13.928\,.
\end{equation}

\textbf{Application to the quark seed.} The SQCD realisation of the
O'Raifeartaigh mechanism identifies the light fermion eigenvalue with
the confined strange quark: $m_- = m_s$. Then the heavy eigenvalue gives
\begin{equation}
m_c^{\mathrm{pred}} = R_{\mathrm{OR}} \times m_s
= 13.928 \times 93.4\;\mathrm{MeV}
= 1301\;\mathrm{MeV}\,.
\end{equation}

\begin{center}
\begin{tabular}{lc}
\toprule
Quantity & Value \\
\midrule
$m_c$ (predicted) & 1301 MeV \\
$m_c$ (PDG 2024) & $1275 \pm 25$ MeV \\
Deviation & $+2.0\%\;\;(+1.0\sigma)$ \\
\bottomrule
\end{tabular}
\end{center}

The prediction uses no fitted parameters beyond the single input
$m_s = 93.4$ MeV ($\overline{\mathrm{MS}}$, 2 GeV).

\subsection{The bion mass relation}

The K\"ahler potential receives non-perturbative corrections from
bion configurations---monopole--instanton pairs in the semiclassical
analysis on $\mathbb{R}^3 \times S^1$~\cite{Unsal2008}. Each
monopole-instanton carries
a single fermionic zero mode per flavour; when lifted by the quark
mass $m_k$, the amplitude acquires a factor $\sqrt{m_k}$.
Since the $N_c = 3$ monopole contributions enter additively in the
effective superpotential, the leading bion correction to the
K\"ahler potential has the form
\begin{equation}
K_{\mathrm{bion}}
= \frac{\zeta^2}{\Lambda^2}
\left|\sum_{k=1}^{N_c} s_k\,\sqrt{\hat{m}_k}\right|^2,
\end{equation}
where $\zeta \propto e^{-S_0/(2N_c)}$ is the monopole fugacity,
$\hat{m}_k$ are the current quark masses, and $s_k = \pm 1$ are signs
determined by the monopole zero-mode structure.

Organising the bion amplitudes into seed-sector and bloom-sector
channels:
\begin{equation}
S_{\mathrm{seed}} = \sqrt{m_s} + \sqrt{m_c}\,,
\qquad
S_{\mathrm{bloom}} = -\sqrt{m_s} + \sqrt{m_c} + \sqrt{m_b}\,,
\end{equation}
the K\"ahler correction contains a term
\begin{equation}
\delta K
\;\sim\;
\lambda_1\,|S_{\mathrm{seed}}|^2
+ \lambda_2\,|S_{\mathrm{bloom}} - 2\,S_{\mathrm{seed}}|^2\,.
\end{equation}
The second factor is a perfect square that vanishes at the minimum.
Setting $S_{\mathrm{bloom}} = 2\,S_{\mathrm{seed}}$ gives the
$v_0$-doubling condition:
\begin{equation}
\boxed{\sqrt{m_b}
= N_c\,\sqrt{m_s} + \sqrt{m_c}
= 3\sqrt{m_s} + \sqrt{m_c}\,.}
\end{equation}

The integer coefficient is the number of colours.

\textbf{Evaluation.} Using the predicted $m_c = 1301$ MeV:
\begin{equation}
\sqrt{m_b}
= 3 \times \sqrt{93.4} + \sqrt{1301}
= 3 \times 9.665 + 36.069
= 65.064\;\mathrm{MeV}^{1/2}\,,
\end{equation}
\begin{equation}
m_b^{\mathrm{pred}} = 4233\;\mathrm{MeV}\,.
\end{equation}

If instead the PDG value $m_c = 1275$ MeV is used:
$m_b = 4186\;\mathrm{MeV}$ ($+0.2\sigma$ from PDG).

\begin{center}
\begin{tabular}{lccc}
\toprule
Input $m_c$ & $m_b$ (predicted) & PDG ($4180 \pm 30$) & Pull \\
\midrule
1301 MeV (from step 1) & 4233 MeV & $+1.8\sigma$ & \\
1275 MeV (PDG) & 4186 MeV & $+0.2\sigma$ & \\
\bottomrule
\end{tabular}
\end{center}

The tighter prediction ($+0.15\%$, $0.2\sigma$) uses the measured
charm mass and tests the bion relation in isolation. The wider one
($+1.3\%$, $1.8\sigma$) propagates the O'Raifeartaigh prediction
and tests the chain.

The principal caveat is that this analysis
applies on $\mathbb{R}^3 \times S^1$; connecting to $\mathbb{R}^4$
physics requires \"Unsal's continuity conjecture~\cite{Unsal2012},
which posits that center-symmetric confinement at small $S^1$ is
continuously connected to the large-$S^1$ confining vacuum.

\subsection{The Yukawa eigenvalue constraint}

The heavy quarks $(c, b, t)$ couple to the electroweak sector through
the Yukawa superpotential. At $\tan\beta = 1$ (required by the mixed
up/down-type structure of the mass triples), the eigenvalue sum rule of
the $3 \times 3$ mass matrix in the $(c, b, t)$ sector imposes:
\begin{equation}
\boxed{\langle z_k^2 \rangle = z_0^2
\quad\text{for}\quad
\zch_k = \sqrt{m_k},\;\; k \in \{c, b, t\}\,.}
\end{equation}

This is the same structural identity that appears at the
O'Raifeartaigh vacuum---an algebraic consequence of the
superpotential, not a fitted parameter.

\textbf{Solving for $m_t$.} Let $x = \sqrt{m_t}$,
$A = \sqrt{m_c} + \sqrt{m_b}$, $\Sigma = m_c + m_b$. The constraint
$3(\Sigma + x^2) = 2(A + x)^2$ gives:
\begin{equation}
\sqrt{m_t}
= 2(\sqrt{m_c} + \sqrt{m_b})
+ \sqrt{3}\,\sqrt{m_c + m_b + 4\sqrt{m_c\,m_b}}\,.
\end{equation}

\textbf{From the self-consistent chain} ($m_c = 1301$, $m_b = 4233$ MeV):
\begin{align}
A &= \sqrt{1301} + \sqrt{4233} = 36.07 + 65.06 = 101.13\,, \notag\\
\sqrt{m_t} &= 2 \times 101.13 + \sqrt{44762} = 202.26 + 211.57 = 413.83\,, \notag\\
m_t^{\mathrm{pred}} &= 171{,}250\;\mathrm{MeV}\,.
\end{align}

\begin{center}
\begin{tabular}{lc}
\toprule
Quantity & Value \\
\midrule
$m_t$ (predicted, chain) & 171,250 MeV \\
$m_t$ (PDG pole) & $172{,}760 \pm 300$ MeV \\
Deviation & $-0.87\%$ \\
\bottomrule
\end{tabular}
\end{center}

\textbf{Scheme dependence.} The chain input $m_s = 93.4\MeV$
is $\overline{\mathrm{MS}}$ at $2\GeV$, while the PDG top mass
$172.57 \pm 0.29\GeV$ is the pole mass. The two differ by the
QCD self-energy: $m_t^{\text{pole}}/\bar m_t(\bar m_t) \approx 1.06$
at two loops, giving $\bar m_t(\bar m_t) = 162.5 \pm 1.1\GeV$.
Converting the prediction to $\overline{\mathrm{MS}}$ at the same
order yields $\bar m_t(\bar m_t) \approx 161.3\GeV$
($-1.1\sigma$ from PDG). The dominant systematic is
the chain's mixed-scale nature---light quarks at $2\GeV$,
heavy quarks effectively at $m_q(m_q)$---which introduces
an $\mathcal{O}(1\GeV)$ scheme ambiguity. The prediction is
therefore stated as $m_t = 171.3 \pm 1.5\GeV$, where the
uncertainty reflects this scheme mismatch.

\subsection{The complete $m_s$ chain}

Three heavy quark masses from one input:
\begin{equation}
\boxed{
\begin{aligned}
m_s &= 93.4\;\mathrm{MeV} \quad\text{(input)}\\[4pt]
&\xrightarrow{\;R_{\mathrm{OR}} = 7 + 4\sqrt{3}\;}
m_c = 1301\;\mathrm{MeV} \quad[+2.0\%]\\[4pt]
&\xrightarrow{\;\sqrt{m_b} = 3\sqrt{m_s} + \sqrt{m_c}\;}
m_b = 4233\;\mathrm{MeV} \quad[+1.3\%]\\[4pt]
&\xrightarrow{\;\langle z_k^2\rangle = z_0^2\text{ on }(c,b,t)\;}
m_t = 171{,}250\;\mathrm{MeV} \quad[-0.9\%]
\end{aligned}
}
\end{equation}

The chain spans a factor of 1834 in mass (from $m_s$ to $m_t$),
controlled by three structural constants of the Lagrangian: the
O'Raifeartaigh mass ratio $R_{\mathrm{OR}}$, the bion doubling
(with coefficient $N_c = 3$), and the eigenvalue sum rule
$\langle z_k^2\rangle = z_0^2$. None of these is fitted; each is an algebraic
consequence of the superpotential at its stabilised vacuum.


%======================================================================
\section{The lepton sector}
\label{sec:leptons}
%======================================================================

\subsection{SU(2) SQCD with $N_f = 3$}

The lepton sector is described by a \emph{separate} confining gauge
theory: $SU(2)$ SQCD with $N_f = 3$ fundamental flavors. This is an
$s$-confining theory ($N_f = N_c + 1 = 3$), meaning the infrared
dynamics is entirely captured by gauge-invariant composites with a
smooth confining superpotential (no chiral symmetry breaking, no runaway).

The UV fields are three chiral superfields in the fundamental of SU(2):
$\Psi^i_a$ with $i = 1,2,3$ (flavor) and $a = 1,2$ (color). The
confined degrees of freedom are the antisymmetric composites
$L^{ij} = \epsilon^{ab}\Psi^i_a\Psi^j_b$ ($i < j$), giving three
independent chiral superfields $L_1 = L^{12}$, $L_2 = L^{13}$,
$L_3 = L^{23}$, each a gauge singlet under $SU(2)$.

The dynamical superpotential from $s$-confinement is the Pfaffian:
$W_{\text{dyn}} = L_1 L_2 L_3 / \Lambda_L^3$.

\subsection{O'Raifeartaigh deformation}

Adding a singlet $X_L$ and the standard O'Raifeartaigh couplings gives
\begin{equation}
W_{\text{lepton}} = f\, X_L + m\, L_1 L_3 + g\, X_L L_1^2
  + \frac{1}{\Lambda_L^3}\, L_1 L_2 L_3\,.
\end{equation}

The metastable SUSY-breaking vacuum sits at
$\langle L_1\rangle = \langle L_3\rangle = 0$,
$\langle X_L\rangle = v$, with $F_{X_L} = f \neq 0$.
The fermion mass matrix at this vacuum has the $2\times 2$ block
\begin{equation}
M_{2\times 2} = \begin{pmatrix} 2gv & m \\ m & 0 \end{pmatrix}
\end{equation}
in the $(L_1, L_3)$ subspace, with eigenvalues
$m_\pm = \sqrt{g^2v^2 + m^2} \pm gv$.
At $gv/m = \sqrt{3}$, the spectrum is
$(0,\; 0,\; (2-\sqrt{3})m,\; (2+\sqrt{3})m)$,
where the two zero eigenvalues correspond to the goldstino ($\psi_{X_L}$)
and the pseudo-modulus fermion ($\psi_{L_2}$).

The K\"ahler potential for the lepton sector is
\begin{equation}
K_{\text{lepton}} = |X_L|^2\!\left(1 - \frac{|X_L|^2}{3\mu_K^2}\right)
  + |L_1|^2 + |L_2|^2 + |L_3|^2\,,
\end{equation}
with the pole condition $\mu_K = \sqrt{2}\,m/g$.

The $s$-confining term $L_1 L_2 L_3/\Lambda_L^3$ has all second
derivatives proportional to fields, hence contributes nothing to the
mass matrices at $\langle L_k\rangle = 0$. The spectrum is therefore
precisely the SU(3) O'Raifeartaigh spectrum plus one extra massless
chiral multiplet ($L_2$), which is a pseudo-modulus stabilised by the
CW potential with a mass suppressed by $(m/\Lambda_L)^6$.
The scalar $B$-term structure ($\pm 2gf$ splitting between
$\sigma$ and $\pi$ sectors) and the identity
$\mathrm{STr}[\mathcal{M}^2] = 0$ carry over identically.

\subsection{Lepton mass prediction}

The O'Raifeartaigh seed for the lepton sector has the spectrum
$(0,\, (2-\sqrt{3})m_{\text{OR}},\, (2+\sqrt{3})m_{\text{OR}})$
with $m_{\text{OR}} = (m_e + m_\mu + m_\tau)/4 = 470.76\MeV$.
The mass-charge identity $\langle z_k^2\rangle = z_0^2$ is satisfied exactly at the seed.

Because the SU(2) theory has no magnetic bion mechanism (the
monopole-instanton structure of SU(2) does not produce the
$v_0$-doubling effect), the physical spectrum stays close to the
O'Raifeartaigh seed. The bloom is a tiny $2.3^\circ$ rotation in
the angular parameter $\delta$, at nearly fixed $v_0$. The
mass-charge identity applied to the charged lepton masses gives:
\begin{equation}
\langle z_k^2 \rangle = z_0^2
\quad\text{for}\quad
\zch_k = \sqrt{m_k},\;\; k \in \{e, \mu, \tau\}\,.
\end{equation}

Solving for $m_\tau$ from the measured $m_e$ and $m_\mu$:
\begin{equation}
m_\tau^{\mathrm{pred}} = 1776.97\;\mathrm{MeV}\,.
\end{equation}

\begin{center}
\begin{tabular}{lc}
\toprule
Quantity & Value \\
\midrule
$m_\tau$ (predicted) & 1776.97 MeV \\
$m_\tau$ (PDG 2024) & $1776.86 \pm 0.12$ MeV \\
Deviation & $+0.006\%\;\;(+0.9\sigma)$ \\
\bottomrule
\end{tabular}
\end{center}

This is the most precise prediction in the programme. Unlike the quark
triples, the lepton triple uses unambiguous pole masses and requires no
$\overline{\mathrm{MS}}$ conventions.

\subsection{Parameter counting for the lepton sector}

Two parameters determine three lepton masses:
\begin{center}
\begin{tabular}{llc}
\toprule
Parameter & What it determines & Count \\
\midrule
$m_{\text{OR}} = 470.76\MeV$ & Overall scale: $M_0 = 4\,m_{\text{OR}}$ & 1 \\
$gv/m = \sqrt{3}$ & Energy-balance condition (constraint) & 0 \\
$\delta_{\text{bloom}} = -2.27^\circ$ & Electron mass generation & 1 \\
\bottomrule
\end{tabular}
\end{center}
Equivalently: specifying any two of $(m_e, m_\mu, m_\tau)$ determines
the third, via $\langle z_k^2\rangle = z_0^2$.


%======================================================================
\section{Electroweak predictions}
\label{sec:electroweak}
%======================================================================

\subsection{The Higgs mass at $\tan\beta = 1$}

At $\tan\beta = 1$, the $D$-term quartic
$(g^2+g'^2)(|H_u|^2 - |H_d|^2)^2/8$ vanishes on the
$D$-flat direction $H_u = H_d$. The tree-level lightest
CP-even Higgs mass is $m_h = m_Z|\cos 2\beta| = 0$---a known
MSSM result.

In the NMSSM, a singlet coupling
$\lambda_S\,S\,H_u\!\cdot\!H_d$ generates an extra quartic
$\lambda_S^2|H_u\!\cdot\!H_d|^2$; at $\sin 2\beta = 1$,
$m_h^2 = \lambda_S^2 v^2/2$, and
$\lambda_S = 0.72$ gives $m_h = 125\GeV$ at tree level.

The Seiberg Lagrange multiplier~$X$ cannot serve as the NMSSM
singlet: $[X] = \text{mass}^{-3}$, so the coupling
$\lambda_0\,X\,H_u\!\cdot\!H_d$ requires
$[\lambda_0] = \text{mass}^4$. Dimensional analysis gives
$\lambda_{\text{NMSSM}} = \hat\lambda$ with
$\hat\lambda \approx 0.51$, but positive meson VEVs
require $\hat\lambda < 2\Lambda^2/v^2 \approx 3\times 10^{-6}$:
a five-order-of-magnitude conflict. Therefore~$S$ must be a
\emph{separate} elementary singlet.

\begin{center}
\begin{tabular}{lc}
\toprule
Quantity & Value \\
\midrule
$m_h$ (tree-level, $\lambda_S = 0.72$) & 125.35 GeV \\
$m_h$ (PDG) & $125.25 \pm 0.17$ GeV \\
Pull & $0.6\sigma$ \\
\bottomrule
\end{tabular}
\end{center}

\subsection{The Cabibbo angle}

The UV quark mass matrix has a nearest-neighbour (Fritzsch) texture
\begin{equation}\label{eq:fritzsch}
M = \begin{pmatrix} 0 & A & 0 \\ A^* & 0 & B \\ 0 & B^* & C \end{pmatrix}\,,
\end{equation}
a structural property of the seesaw vacuum $M \propto \hat{m}^{-T}$
that inverts the mass hierarchy while preserving eigenvectors.
In the hierarchy $|A| \ll |B| \ll C$, the eigenvalues are
$m_1 \approx |A|^2 C/|B|^2$, $m_2 \approx |B|^2/C$, $m_3 \approx C$,
and the $1$--$2$ mixing angle satisfies the
Gatto--Sartori--Tonin relation~\cite{GST1968}:
\begin{equation}
\sin\theta_C = \sqrt{\frac{m_d}{m_s}}
= \sqrt{\frac{4.67}{93.4}}
= 0.2236\,.
\end{equation}

\paragraph{Up-type correction.}
With separate Fritzsch textures for up-type and down-type quarks,
the CKM element receives a phase-dependent correction:
\begin{equation}\label{eq:fritzsch-Vus}
|V_{us}| = \left|\sqrt{\frac{m_d}{m_s}}
  - e^{i\phi}\sqrt{\frac{m_u}{m_c}}\right|\,.
\end{equation}
The second term $\sqrt{m_u/m_c} = 0.041$ is an 18\% correction.
With $\phi = \pi/2$ the classic Fritzsch prediction gives
$|V_{us}| = 0.2274$ ($3.8\sigma$ high); the PDG value
$0.2243 \pm 0.0008$ is reproduced at $\phi \approx 86^\circ$.

\begin{center}
\begin{tabular}{lcc}
\toprule
Method & $\sin\theta_C$ & $\theta_C$ \\
\midrule
$\sqrt{m_d/m_s}$ (GST) & 0.2236 & $12.92^\circ$ \\
$|V_{us}|$ (PDG 2024) & $0.2243 \pm 0.0008$ & $12.96^\circ$ \\
Deviation & $-0.3\%$ & $(-0.9\sigma)$ \\
\bottomrule
\end{tabular}
\end{center}

\paragraph{Charge-language interpretation.}
In the angular parametrisation~\eqref{eq:zk-angular}
with $k = 0,1,2$ labelling $(d, s, b)$, the GST ratio becomes
\begin{equation}
\sin\theta_C
= \frac{|1 + \sqrt{2}\cos\delta|}
       {|1 + \sqrt{2}\cos(\delta + 2\pi/3)|}\,.
\end{equation}
At the seed $\delta = 3\pi/4$, the numerator vanishes
($1 + \sqrt{2}\cos(3\pi/4) = 0$), forcing $m_d = 0$ and
$\theta_C = 0$.
A bloom $\delta = 3\pi/4 + \varepsilon$ opens both a nonzero
$m_d$ and a nonzero Cabibbo angle; they are the \emph{same degree
of freedom}.  The mass-charge identity
$\langle z_k^2\rangle = z_0^2$ is an identity
$\sum \cos^2(\delta + 2\pi k/3) = 3/2$ for all~$\delta$ and
imposes no constraint on~$\varepsilon$.  The Cabibbo angle
is thus controlled entirely by~$\delta$: it measures how far
the down-type charge configuration has bloomed from the massless
seed.  The scale parameter $z_0$ remains separate.

This is not an independent prediction of the superpotential
but a structural property of the seesaw-inverted mass hierarchy.
The free parameters $m_u$ and $m_d$ sit outside the mass chain;
the GST relation connects them to the CKM matrix.

\subsection{The Casimir eigenvalue equation}

Consider the eigenvalue equation
\begin{equation}
x^2 \;+\; C_2\,x \;-\; C_2 = 0\,,
\qquad C_2 = s(s+1)\,,
\end{equation}
where $C_2$ is the quadratic Casimir of $SU(2)$ at spin $s$.
Three structural constraints pin the polynomial: (i) no spin-independent
mass term, (ii) a massless eigenstate at $s = 0$,
(iii) coefficient antisymmetry. The solutions are
\begin{equation}
x_{\pm}(s) = \frac{-C_2 \pm \sqrt{C_2^2 + 4C_2}}{2}\,.
\end{equation}

Defining the ratio
$R \equiv 1 - x_+(1/2)/x_+(1)$ and evaluating at
$s = 1/2$ ($C_2 = 3/4$) and $s = 1$ ($C_2 = 2$):
\begin{equation}
R = \frac{(\sqrt{19} - 3)(\sqrt{19} - \sqrt{3}\,)}{16}
  \;\approx\; 0.2231013\ldots
\end{equation}
This should be compared with the on-shell electroweak mixing
parameter $\sin^2\!\theta_W^{\text{os}} = 0.22320 \pm 0.00026$
(PDG 2024).

The $W$-mass prediction from $M_W = M_Z\sqrt{1 - R}$:
\begin{equation}
M_W^{\text{pred}}
  = 91.1876 \times \sqrt{1 - 0.22310}
  = 80.374 \pm 0.002\GeV\,,
\end{equation}
which lies $0.39\sigma$ from the PDG average
$M_W = 80.369 \pm 0.013\GeV$ and strongly disfavors the
anomalous CDF-II measurement ($6.2\sigma$ tension).

The dynamical origin of this equation and the choice of
renormalization scheme (on-shell vs.\ $\overline{\text{MS}}$)
remain open problems.


%======================================================================
\section{Complete particle spectrum}
\label{sec:spectrum}
%======================================================================

\subsection{O'Raifeartaigh model spectrum}

At general $t = gv/m$, the scalar potential $V = \sum|F_k|^2$
depends on a second dimensionless parameter $y = gf/m^2$ (the
SUSY-breaking strength, with $y < 1/2$ for stability).
The $6\times 6$ real scalar mass matrix in the basis
$(\sigma_0,\sigma_1,\sigma_2,\pi_0,\pi_1,\pi_2)$ is block-diagonal,
$\mathcal{M}^2_S = \mathrm{diag}(\mathcal{M}^2_\sigma,
\mathcal{M}^2_\pi)$, with
\begin{equation}
  \frac{\mathcal{M}^2_\sigma}{m^2}
  = \begin{pmatrix} 0 & 0 & 0 \\ 0 & 4t^2{+}2y{+}1 & 2t \\
    0 & 2t & 1 \end{pmatrix}, \quad
  \frac{\mathcal{M}^2_\pi}{m^2}
  = \begin{pmatrix} 0 & 0 & 0 \\ 0 & 4t^2{-}2y{+}1 & 2t \\
    0 & 2t & 1 \end{pmatrix}.
\end{equation}
The $B$-term $B_{11} = 2gf$ is the sole nonvanishing element of
$B_{ij} = \sum_k F_k W_{kij}$, contributing $+2y$ to $\sigma$
and $-2y$ to $\pi$. The four nonzero eigenvalues are
\begin{equation}
  \frac{m^2_{\sigma,\pm}}{m^2}
    = 2t^2 + y + 1 \pm \sqrt{(2t^2+y)^2+4t^2}\,,\quad
  \frac{m^2_{\pi,\pm}}{m^2}
    = 2t^2 - y + 1 \pm \sqrt{(2t^2-y)^2+4t^2}\,.
\end{equation}
At $t = \sqrt{3}$, each fermion mass is bracketed by its
$\sigma$ (heavier) and $\pi$ (lighter) scalar partners.
In the SUSY limit $y \to 0$, the $\sigma$--$\pi$ splittings vanish
and each fermion mass-squared $m^2_\pm = 7 \pm 4\sqrt{3}$ is
doubly degenerate.

The supertrace vanishes identically:
$\mathrm{STr}[\mathcal{M}^2] = m^2(8t^2{+}4) - 2m^2(4t^2{+}2) = 0$
for all $t$ and $y$, because the $B$-term contributions
$\pm 2y$ cancel between sectors.

\subsection{Scalar meson masses}

At the Seiberg vacuum with soft breaking $\tilde{m}^2 \sim \fpi^2$:
\begin{itemize}
\item 12 off-diagonal real scalar modes: all degenerate at
$m = \sqrt{2}\,\fpi \approx 130\MeV$.
\item 6 diagonal real scalar modes: all at
$m = \fpi = 92\MeV$.
\end{itemize}
All scalar meson masses are positive. The
Yukawa-induced $F$-terms are nonzero at the seesaw vacuum but do not
couple to the off-diagonal meson directions: the off-diagonal
cofactors vanish at a diagonal matrix. A fundamental block-diagonal
obstruction prevents cross-flavour meson mixing.

\subsection{Mesino fermion masses}

The off-diagonal mesinos form three Dirac pairs with masses set by
the inverted seesaw hierarchy $|X_0|\langle M_k\rangle \propto 1/m_k$:

\begin{center}
\begin{tabular}{lcc}
\toprule
Pair & Spectator & Dirac mass \\
\midrule
$\psi(M^u_d, M^d_u)$ & $s$ & $1.1 \times 10^{-5}\MeV$ \\
$\psi(M^u_s, M^s_u)$ & $d$ & $2.3 \times 10^{-4}\MeV$ \\
$\psi(M^d_s, M^s_d)$ & $u$ & $4.9 \times 10^{-4}\MeV$ \\
\bottomrule
\end{tabular}
\end{center}

The mass hierarchy is inverted relative to quark masses: off-diagonal
mesinos involving the lightest quarks are heaviest. The ISS
Coleman--Weinberg potential from these loops transmits the quark mass
ratios through the confined spectrum.

\subsection{Baryons}

With $m_B = \Lambda = 300\MeV$ in the superpotential
$W_B = m_B\,B\tilde{B}$:
scalar baryons at $300\MeV$ and a Dirac baryonino pair at $300\MeV$.
SUSY is approximately preserved in the baryon sector.

\subsection{Heavy quark sector}

\begin{center}
\begin{tabular}{lccc}
\toprule
Quark & $Q_{\text{em}}$ & Mass & Origin \\
\midrule
charm & $+2/3$ & 1275 MeV (1301 predicted) & O'Raifeartaigh seed \\
bottom & $-1/3$ & 4180 MeV (4233 predicted) & Bion mass relation \\
top & $+2/3$ & 172,760 MeV & Direct Yukawa $y_t$ \\
\bottomrule
\end{tabular}
\end{center}

\subsection{Higgs sector}

The lightest CP-even Higgs at $m_h = \lambda_S v/\sqrt{2} = 125.35\GeV$
($0.6\sigma$ from PDG). Heavy NMSSM states ($H$, $A$, $H^\pm$, singlino)
depend on soft parameters not specified by the confined sector.

\subsection{Supertrace}

$\text{STr}[M^2] = 2\operatorname{Tr}(M^2_{\text{soft}})
= 18\,\fpi^2 = 152{,}352\MeV^2$:
the $W_{IJ}^2$ piece cancels between scalar and fermion sectors,
and the holomorphic third-derivative piece cancels between real
and imaginary scalar components, leaving only the nine
soft meson masses.

\subsection{Parameter budget}

\begin{center}
\small
\begin{tabular}{p{5cm}p{3.5cm}p{5.5cm}}
\toprule
Relation & Predicts & From \\
\midrule
O'Raifeartaigh: $m_c/m_s = (2{+}\sqrt{3})^2$ & $m_c$ from $m_s$ & $W$: $gv/m = \sqrt{3}$ \\
Bion: $\sqrt{m_b} = 3\sqrt{m_s} + \sqrt{m_c}$ & $m_b$ from $m_s, m_c$ & $K_{\text{bion}}$: minimum \\
Yukawa: $\langle z_k^2\rangle = z_0^2$ on $(c,b,t)$ & $m_t$ from $m_c, m_b$ & $W_{\text{Yukawa}}$ at $\tan\beta = 1$ \\
Lepton: $\langle z_k^2\rangle = z_0^2$ on $(e,\mu,\tau)$ & $m_\tau$ from $m_e, m_\mu$ & $W_{\text{lepton}}$: O'Raifeartaigh \\
GST: $\sin\theta_C = \sqrt{m_d/m_s}$ & $\theta_{12}$ from $m_d, m_s$ & Seesaw structure \\
\midrule
\textbf{Remaining free:} & \multicolumn{2}{l}{$m_u,\, m_d,\, m_s,\, m_e,\, m_\mu,\, \theta_{23},\, \theta_{13},\, \delta_{\rm CP}$ (8 parameters)} \\
\bottomrule
\end{tabular}
\end{center}


%======================================================================
\section{Summary of predictions}
\label{sec:summary_pred}
%======================================================================

\begin{table}[ht]
\centering
\small
\begin{tabular}{llll}
\toprule
Quantity & Predicted & PDG 2024 & Source \\
\midrule
$m_c$ & $1{,}301\MeV$ & $1{,}275 \pm 25$ ($1.0\sigma$)
  & O'Raifeartaigh: $gv/m = \sqrt{3}$ \\
$m_b$ & $4{,}233\MeV$ & $4{,}180 \pm 30$ ($1.8\sigma$)
  & Bion K\"ahler \\
$m_t$ & $171.3 \pm 1.5\GeV$ & $172.76 \pm 0.30$
  & Yukawa eigenvalue constraint \\
$m_\tau$ & $1776.97\MeV$ & $1776.86 \pm 0.12$ ($0.9\sigma$)
  & Lepton mass-charge identity \\
$\sin^2\theta_W$ & $0.22310$ & $0.22320 \pm 0.00026$ ($0.4\sigma$)
  & Casimir equation \\
$V_{us}$ & $0.224$ & $0.2250 \pm 0.0007$ ($2.0\sigma$)
  & GST relation \\
$m_h$ & $125.35\GeV$ & $125.25 \pm 0.17$ ($0.6\sigma$)
  & NMSSM: $\lambda_S v/\sqrt{2}$ \\
$\tan\beta$ & $1$ & --- (not directly measured)
  & Mixed-triple structure \\
$\text{STr}[M^2]$ & $18\,\fpi^2$ & --- (not directly measured)
  & Soft mass only \\
\bottomrule
\end{tabular}
\caption{Predictions from the sBootstrap Lagrangian.
The quark masses use $m_s = 93.4\MeV$ as input;
the bion $m_b$ propagates the O'Raifeartaigh-predicted $m_c$.
The lepton prediction uses $m_e$ and $m_\mu$ as inputs.
The Casimir and GST predictions are independent of~$m_s$.}
\label{tab:predictions}
\end{table}


%======================================================================
\section{Open problems}
\label{sec:open}
%======================================================================

\subsection{The diquark sector}
The Pauli theorem (Section~\ref{sec:diquarks}) proves that no $J=0$,
color-$\bar{3}$ diquark can have symmetric flavor. The sBootstrap's
$\mathbf{15}$ diquark must therefore be non-composite:
either an algebraic image of the Miyazawa supercharge or a fundamental
$SO(32)$ string state. However, the mixed anticommutator
$\{Q_\alpha, \mathcal{Q}_{[ij],\beta}\} = \epsilon_{\alpha\beta}\,
\mathcal{S}_{[ij]}$ produces a diquark generator in
$\overline{\mathbf{10}}$ (antisymmetric), not the symmetric
$\mathbf{15}$; thus the Miyazawa route alone does not resolve the
conflict. The exotic $+4/3$-charged
diquarks of the $\mathbf{15}$ would be a testable consequence.

\subsection{The bloom mechanism}

The mass-charge identity constrains but does not uniquely fix the
angular parameter $\delta$ that controls how the seed develops into a
full triple. The O'Raifeartaigh vacuum fixes the seed ratio
$gv/m = \sqrt{3}$ but leaves $\delta$ undetermined.

The lowest-dimension $R$-breaking deformation that affects the
spectrum is the bilinear $\Delta W = \epsilon\,\Phi_0\Phi_2$, which
forces $\langle\Phi_1\rangle \neq 0$ and mixes the goldstino with the
massive sector.  (Cubic deformations in $\Phi_1,\Phi_2$ alone vanish
at the vacuum and are inert.)  To first order in $\varepsilon =
\epsilon/m$, all eigenvalue shifts vanish, but the bloom angle shifts
at $O(\varepsilon)$ because $\delta\delta \sim \sqrt{m_0} \sim
\varepsilon$:
$$
  \delta\delta = 3^{5/4}\,|\varepsilon| \approx 3.95\,|\varepsilon|\,.
$$
The direction is \emph{toward} the physical quark spectrum
($\delta > 3\pi/4$).  The determinant condition $m_+m_- = m^2$ is
broken at $O(\varepsilon^2)$ with fractional deviation
$49\,\varepsilon^2$.  The seed is therefore an unstable fixed point
under any $R$-breaking, and the instability points in the right
direction. The bloom magnitude, however, requires a nonperturbative
mechanism.

A dynamical origin for the lepton phase exists in SQCD: for $N_f = N_c = 3$,
instantons break $U(1)_{R+2A} \to \mathbb{Z}_{18}$, which contains
$\mathbb{Z}_9$ as a subgroup. The three-instanton operator
$(\det M)^3/\Lambda^{18}$ generates a $\cos(9\delta)$ potential whose
minima include the observed lepton mass-spectrum phase to $0.9\sigma$
in~$m_\tau$.

For the quark sector, the combined potential of the frozen-sign bion
(which opposes the bloom) and the three-instanton term (which drives
the bloom) has a stable equilibrium within $2.8^\circ$ of the physical
quark bloom position.

\subsection{Lepton masses are not mesino masses}

The fermion mass matrix $W_{IJ}$ at the Seiberg vacuum has
light eigenvalues ($\sim 10^{-3}$--$10^{-2}\MeV$) too small
and too degenerate ($\langle z_k^2\rangle/z_0^2 = 0.44$, far from unity) to reproduce lepton masses.
The obstruction is structural: K\"ahler rescaling is a congruence
transformation that preserves the rank of~$W_{IJ}$ and cannot generate
new mass eigenvalues. Lepton masses therefore require the separate
SU(2) confining sector described in Section~\ref{sec:leptons}.

\subsection{The $SO(32)$ embedding}

The $SO(32)$ gauge group of Type~I / heterotic string theory
admits the maximal-subgroup chain
$SO(32) \supset SU(16) \times U(1)_A
  \supset SU(5) \times SU(3) \times U(1)_A \times U(1)_B$,
where $\mathbf{16} = (\mathbf{5},\mathbf{3})_{+1}
  \oplus (\mathbf{1},\mathbf{1})_{-15}$
embeds the five quark flavors times three colors as a single
tensor-product block.

\paragraph{Complete decomposition.}
The adjoint $\mathbf{496} = \wedge^2(\mathbf{32})$ splits as
$(\mathbf{255}\!\oplus\!\mathbf{1})_{0}
  \oplus \mathbf{120}_{+2}
  \oplus \overline{\mathbf{120}}_{-2}$
under $SU(16) \times U(1)_A$.  The $SU(5) \times SU(3)$ content:

\begin{center}
\begin{tabular}{llcccl}
\toprule
Sector & $SU(5)\times SU(3)$ & dim & $q_A$ & $q_B$ & Origin \\
\midrule
$\mathbf{255}\!\oplus\!\mathbf{1}$
 & $(\mathbf{24},\mathbf{8})$ & 192 & 0 & 0
   & $\mathrm{adj}(SU(15))$ \\
 & $\boldsymbol{(24,1)}$ & $\boldsymbol{24}$ & 0 & 0
   & $\mathrm{adj}(SU(15))$ \\
 & $(\mathbf{1},\mathbf{8})$ & 8 & 0 & 0
   & $\mathrm{adj}(SU(15))$ \\
 & $(\mathbf{5},\mathbf{3})\!+\!(\bar{\mathbf{5}},\bar{\mathbf{3}})$
   & 30 & 0 & $\pm 16$ & cross-terms \\
 & $2\times(\mathbf{1},\mathbf{1})$ & 2 & 0 & 0
   & $U(1)_{A,B}$ \\
\midrule
$\mathbf{120}$
 & $\boldsymbol{(15,\bar{3})}$ & $\boldsymbol{45}$ & $+2$ & $+2$
   & $S^2(\mathbf{5})\otimes\wedge^2(\mathbf{3})$ \\
 & $(\overline{\mathbf{10}},\mathbf{6})$ & 60 & $+2$ & $+2$
   & $\wedge^2(\mathbf{5})\otimes S^2(\mathbf{3})$ \\
 & $(\mathbf{5},\mathbf{3})$ & 15 & $+2$ & $-14$
   & cross-term \\
\midrule
$\overline{\mathbf{120}}$
 & $\boldsymbol{(\overline{15},3)}$ & $\boldsymbol{45}$ & $-2$ & $-2$
   & conjugate \\
 & $(\mathbf{10},\bar{\mathbf{6}})$ & 60 & $-2$ & $-2$
   & conjugate \\
 & $(\bar{\mathbf{5}},\bar{\mathbf{3}})$ & 15 & $-2$ & $+14$
   & conjugate \\
\midrule
 & \textbf{Total} & \textbf{496} & & & \\
\bottomrule
\end{tabular}
\end{center}

\noindent
The target---$(\mathbf{24},\mathbf{1})$,
$(\mathbf{15},\bar{\mathbf{3}})$,
$(\overline{\mathbf{15}},\mathbf{3})$---accounts
for $24 + 45 + 45 = 114$ of the 496 states, each appearing once.

\paragraph{Why the symmetric $\mathbf{15}$ appears.}
The nontrivial content of $\wedge^2(\mathbf{16})$ comes from
$\wedge^2[(\mathbf{5},\mathbf{3})]$.  The key identity is
$$
\wedge^2(V \otimes W)
  \;=\; \bigl[S^2(V) \otimes \wedge^2(W)\bigr]
  \;\oplus\; \bigl[\wedge^2(V) \otimes S^2(W)\bigr]\,.
$$
Overall antisymmetry of the composite index $(ia)$ forces a
correlation: symmetric in one factor implies antisymmetric in the
other.  Setting $V = \mathbf{5}$, $W = \mathbf{3}$:
$$
\wedge^2\bigl[(\mathbf{5},\mathbf{3})\bigr]
  \;=\; (\mathbf{15},\bar{\mathbf{3}})
  \;\oplus\; (\overline{\mathbf{10}},\mathbf{6})\,.
$$
The symmetric flavor $\mathbf{15}$ is forced into antisymmetric
color $\bar{\mathbf{3}}$; the pairing is dictated by the
tensor-product structure, not by a modeling choice.

\paragraph{Minimality and obstructions.}
The identity requires $(\mathbf{5},\mathbf{3})$ as a
tensor-product block in the fundamental, hence dimension
$\geq 15$ complex $= 30$ real; $SO(30)$ is the minimal
orthogonal group, and $SO(32)$ adds only two inert singlets.
$SO(16) \times SO(16)$ fails: no single $\mathbf{16}$ can hold
the 15-dimensional block, so $SU(5)$ and $SU(3)$ embed as
separate direct summands whose $\wedge^2$ produces only
$\overline{\mathbf{10}}$ and $\bar{\mathbf{3}}$, never
$\mathbf{15}$.
$E_8 \times E_8$ is likewise blocked: every $E_8$ decomposition
under $SU(5)$ yields only exterior powers
($\mathbf{5}$, $\overline{\mathbf{10}}$, $\mathbf{24}$, \ldots).
The symmetric $\mathbf{15} = S^2(\mathbf{5})$ is algebraically
forbidden.

\paragraph{Orientifold projection.}
Wilson lines and orbifold projections cannot separate
$(\mathbf{15},\bar{\mathbf{3}})$ from
$(\overline{\mathbf{10}},\mathbf{6})$: both carry identical
$U(1)$ charges $(q_A, q_B) = (+2,+2)$, and Schur's lemma forces
any orbifold element commuting with $SU(5) \times SU(3)$ to act
as a scalar on $(\mathbf{5},\mathbf{3})$ and hence on its exterior
power.
A mixed orientifold resolves this.  Placing the $SU(5)$ flavor
branes near an $O^-$ plane (Sp-type, selecting symmetric
$\mathbf{15}$) and the $SU(3)$ color branes near an $O^+$ plane
(SO-type, selecting antisymmetric $\bar{\mathbf{3}}$), worldsheet
parity~$\Omega$ acts with opposite signs on the two tensor
factors and doubly projects out
$(\overline{\mathbf{10}},\mathbf{6})$.  The remaining
$(\mathbf{24},\mathbf{8})$ and bifundamental copies acquire
compactification-scale masses from brane separation; the $U(1)$
generators become massive via the St\"uckelberg mechanism.

\paragraph{A Type~I construction.}
Consider Type~IIB on $T^6/\mathbb{Z}_3$ (each $T^2$ at the $SU(3)$
root lattice) with orientifold projection~$\Omega$.  The $\mathbb{Z}_3$
acts as $(z_1,z_2,z_3) \to (\omega z_1, \omega z_2, \omega z_3)$
with $\omega = e^{2\pi i/3}$, preserving $\mathcal{N}=1$ in 4d.
The Chan-Paton embedding on 16 D9-brane pairs is
$$
  \gamma_\theta = \mathrm{diag}(\omega\cdot\mathbb{1}_5,\;
    \omega^2\cdot\mathbb{1}_5,\; \mathbb{1}_3,\; \mathbb{1}_3)\,.
$$
The orientifold~$\Omega$ swaps the charge-1 and charge-2 blocks
(each of dimension~5), and also swaps the two charge-0 3-blocks.
This gives
$$
  \mathrm{Tr}(\gamma_\theta) = 5\omega + 5\omega^2 + 3 + 3
  = -5 + 6 = 1\,,
$$
satisfying the twisted tadpole without D5-branes.
The constraints $n_1 = n_2$ (orientifold) and
$\mathrm{Tr}(\gamma_\theta) = 1$ uniquely force $n_1 = n_2 = 5$
with 6 charge-0 branes.  Among the 23 possible orientifold actions
on the charge-0 sector, the partition $[3,3]$ with
$\Omega$-swapped blocks is the \emph{unique} embedding yielding
$SU(3)$ as a direct gauge factor.  The resulting gauge group is
$$
  G = SU(5)_{\mathrm{f}} \times SU(3)_{\mathrm{c}}
  \times U(1)^2\,,
$$
with three chiral generations from the $\mathbb{Z}_3$ structure.
The diquark representation (symmetric~$\mathbf{15}$ vs.\
antisymmetric~$\overline{\mathbf{10}}$) is determined by the
orientifold sign on the charge-1/charge-2 pair; the odd
dimension of the flavor block ($n = 5$) rules out the Sp-type
projection, so the symmetric $\mathbf{15}$ is the unique
consistent choice.
The self-referential numerology---$5 \times 3 + 1 = 16$
Chan-Paton factors forming $SO(32)$, whose adjoint contains
exactly the composites of those quarks---awaits a complete
string vacuum.

\subsection{The baryon sector}

The meson--baryon connection requires a Miyazawa
supercharge~\cite{Miyazawa, CattoGursey} in the
$\overline{\mathbf{10}}$ of $SU(5)_f$, with anticommutator
decomposition
$\overline{\mathbf{10}}\otimes\mathbf{10}
  = \mathbf{1}\oplus\mathbf{24}\oplus\mathbf{75}$.
This evades the Haag--\L{}opusza\'nski--Sohnius theorem because
it is an approximate symmetry broken at scale~$\Lambda_{\text{QCD}}$.

\subsection{Vacuum stability and FCNC}

The seesaw vacuum is cosmologically stable:
Coleman--de~Luccia bounce-action estimates give
$S_4 \sim 2 \times 10^7$, far above the threshold
$S_4 > 400$ for a lifetime exceeding the age of the universe.
All off-diagonal meson masses are positive
($m^2 \approx 2\fpi^2$), and a block-diagonal obstruction
prevents cross-flavour condensation. The flavour-universal
Yukawa identity $y_j\langle M_j\rangle = C/v_{\text{EW}}$
automatically satisfies the Glashow--Weinberg condition for the
absence of tree-level Higgs-mediated FCNCs.

\subsection{Generation uniqueness}
The Diophantine system uniquely selects $N = 3$ generations,
but the third equation ($r^2 + s^2 - 1 = 4N$) has the weakest
physical motivation. A dynamical derivation of the counting
constraints would strengthen the argument considerably.


%======================================================================
\section{Conclusions}
\label{sec:conclusions}
%======================================================================

We have constructed an explicit $\mathcal{N}=1$ supersymmetric theory
with gauge group $\mathrm{SU}(3) \times \mathrm{SU}(2)$ and demonstrated
the following principal results.

\paragraph{1.\ The Diophantine counting uniquely selects $N=3$.}
Three equations on the scalar content---$rs = 2N$, $r(r+1)/2 = 2N$,
$r^2+s^2-1 = 4N$---have the unique positive-integer solution
$(N,r,s) = (3,3,2)$, which maps to the $SU(5)$ representations
$\mathbf{24} \oplus \mathbf{15} \oplus \overline{\mathbf{15}}$.
This combination is anomaly-free and asymptotically free ($b_0 = 31/3$).
Right-handed neutrinos are required; the top quark is automatically
excluded from the composite sector.

\paragraph{2.\ The O'Raifeartaigh superpotential produces the
mass-charge identity as a theorem.}
With three chiral superfields, the O'Raifeartaigh model at
$gv/m = \sqrt{3}$ (stabilized by a K\"ahler pole at $c = -1/12$)
produces the fermion spectrum
$(0,\, (2{-}\sqrt{3})\,m,\, (2{+}\sqrt{3})\,m)$
satisfying $\langle z_k^2\rangle = z_0^2$ exactly. This identity is a
boundary condition on the superpotential, not an RG attractor
(one-loop RG drives $\langle z_k^2\rangle/z_0^2 \to 0$,
the opposite extreme).

\paragraph{3.\ The quark mass chain from a single input.}
From $m_s = 93.4\MeV$ alone, the Lagrangian predicts
$m_c = 1301\MeV$ ($+2.0\%$), $m_b = 4233\MeV$ ($+1.3\%$),
and $m_t = 171.3 \pm 1.5\GeV$ ($-0.9\%$), via three structural
constants: the O'Raifeartaigh mass ratio $(2{+}\sqrt{3})^2$, the
bion doubling with coefficient $N_c = 3$, and the Yukawa eigenvalue
constraint $\langle z_k^2\rangle = z_0^2$ on the $(c,b,t)$ charges.

\paragraph{4.\ The lepton sector from SU(2) $s$-confinement.}
A separate $SU(2)$ confining sector with $N_f = 3$ produces the
same O'Raifeartaigh seed spectrum, predicting
$m_\tau = 1776.97\MeV$ ($+0.006\%$, $0.9\sigma$) from
$m_e$ and $m_\mu$.

\paragraph{5.\ Mixed mass triples require $\tan\beta = 1$.}
Since all quark mass triples satisfying $\langle z_k^2\rangle \approx z_0^2$
obligatorily mix up-type and down-type quarks, the mass-charge
identity on masses equals the identity on Yukawa couplings
only at $\tan\beta = 1$. The NMSSM singlet coupling
$\lambda_S = 0.72$ then gives $m_h = 125.35\GeV$ at tree level
($0.6\sigma$ from PDG).

\paragraph{6.\ Five parameters eliminated.}
The O'Raifeartaigh mass ratio, the bion mass relation, the Yukawa
eigenvalue constraint, the lepton mass-charge identity, and
the GST relation together determine five of the Standard Model's
thirteen charged-fermion parameters. The remaining eight are
$m_u$, $m_d$, $m_s$, $m_e$, $m_\mu$, $\theta_{23}$,
$\theta_{13}$, and $\delta_{\rm CP}$, with $m_u$ and $m_d$
connected to the Cabibbo angle.

\paragraph{7.\ The diquark obstruction and SO(32) embedding.}
The symmetric $\mathbf{15}$ required by the counting cannot be
realized as a $qq$ composite (by a kinematic theorem independent of
dynamics). The $SO(32)$ adjoint provides a UV completion containing
both the $\mathbf{24}$ and $\mathbf{15} \oplus \overline{\mathbf{15}}$;
the alternative $E_8 \times E_8$ is blocked. The bloom mechanism,
the $SO(32)$ projection, and a Lagrangian derivation of the
Diophantine counting remain the central open problems.

The empirical evidence for the mass-charge identity---the
precision of the lepton prediction, the accuracy of the meson
triple, and the statistical significance of the quark
chains---will be presented in a companion paper~\cite{companion}.


\begin{thebibliography}{99}

\bibitem{Rivero2005}
A.~Rivero, ``An interpretation of scalars in SO(32),''
Eur.\ Phys.\ J.\ C \textbf{84}, 1058 (2024),
arXiv:2407.05397 [hep-ph]; earlier versions: viXra:0503.001 (2005),
viXra:0809.001 (2008).

\bibitem{companion}
A.~Rivero, ``The mass-charge identity
$\langle z_k^2\rangle = z_0^2$ in fermion mass triples:
empirical evidence,'' in preparation.

\bibitem{Seiberg1995}
N.~Seiberg, ``Electric--magnetic duality in supersymmetric
non-abelian gauge theories,'' Nucl.\ Phys.\ B \textbf{435}, 129
(1995).

\bibitem{ORaifeartaigh}
L.~O'Raifeartaigh, ``Spontaneous symmetry breaking for chiral
scalar superfields,'' Nucl.\ Phys.\ B \textbf{96}, 331 (1975).

\bibitem{ISS2006}
K.~Intriligator, N.~Seiberg, D.~Shih, ``Dynamical SUSY breaking
in meta-stable vacua,'' JHEP \textbf{0604}, 021 (2006),
arXiv:hep-ph/0602239.

\bibitem{Unsal2008}
M.~\"Unsal, ``Abelian duality, confinement, and chiral symmetry
breaking in QCD(adj),'' Phys.\ Rev.\ Lett.\ \textbf{100}, 032005
(2008), arXiv:0708.1772.

\bibitem{Unsal2012}
M.~\"Unsal, ``Theta dependence, sign problems and topological
interference,'' Phys.\ Rev.\ D \textbf{86}, 105012 (2012),
arXiv:1201.6426.

\bibitem{CattoGursey}
S.~Catto and F.~G\"ursey, ``Algebraic treatment of effective
supersymmetry,'' Nuovo Cim.\ A \textbf{86}, 201 (1985).

\bibitem{Miyazawa}
H.~Miyazawa, ``Baryon Number Changing Currents,''
Prog.\ Theor.\ Phys.\ \textbf{36}, 1266 (1966).

\bibitem{GST1968}
R.~Gatto, G.~Sartori, and M.~Tonin, ``Weak self-masses, Cabibbo
angle, and broken $SU(3)$,'' Phys.\ Lett.\ B \textbf{28}, 128 (1968).

\bibitem{VolkovAkulov1973}
D.~V.~Volkov and V.~P.~Akulov, ``Is the neutrino a Goldstone
particle?,'' Phys.\ Lett.\ B \textbf{46}, 109 (1973).

\bibitem{GOR1968}
M.~Gell-Mann, R.~J.~Oakes, B.~Renner, ``Behavior of current
divergences under $SU(3) \times SU(3)$,''
Phys.\ Rev.\ \textbf{175}, 2195 (1968).

\bibitem{FayetIliopoulos}
P.~Fayet and J.~Iliopoulos, ``Spontaneously broken supergauge
symmetries and Goldstone spinors,'' Phys.\ Lett.\ B \textbf{51},
461 (1974).

\bibitem{DGMLY1967}
T.~Das, G.~S.~Guralnik, V.~S.~Mathur, F.~E.~Low, J.~E.~Young,
``Electromagnetic mass difference of pions,''
Phys.\ Rev.\ Lett.\ \textbf{18}, 759 (1967).

\bibitem{Rivero2005b}
A.~Rivero, ``The de Vries formula for $\sin^2\theta_W$,''
arXiv:hep-ph/0510068 (2005).

\bibitem{UTfit}
M.~Bona et~al.\ (UTfit Collaboration), ``Model-independent
constraints on $\Delta F = 2$ operators and the scale of new
physics,'' JHEP \textbf{03}, 049 (2008),
arXiv:0707.0636.

\end{thebibliography}

\end{document}
